\appendixchapter{Tensor Products and Exterior Powers}
\label{app:tensors}
\chaptermark{Tensor Products and Exterior Powers}
\section*{Introduction}
Multilinear rules become easier to organize when they are represented by
linear maps on a single vector space. This appendix develops that advantage
over a field $F$, following Chapter~8 of the author's algebra notes.
The route is construction,
associativity, symmetry, direct sums, finite multilinear products, induced
maps and calculations, Tensor--Hom, exterior powers, and finally the
exterior algebra and alternating forms. A quotient-space preparation makes
the constructions self-contained. No topology or ordering of $F$ is needed;
$F=\R$ connects the final constructions to differential forms.

There are finitely many tensor factors and direct summands. Quotient and
tensor constructions can use arbitrary vector spaces; the formal-symbol
spaces in their proofs are usually infinite dimensional. Dimension formulas
and duality identifications carry their finite-dimensional hypotheses explicitly.
The exterior algebra below is formed for a finite-dimensional space, so
only finitely many of its homogeneous components are nonzero.

The recurring task is to construct a map without choosing a representation
of its input. First define a map on independent vectors, verify its slot
identities, and factor it through the relevant universal property. Construct
an inverse by the same method when an isomorphism is claimed. Once both maps
exist, agreement on spanning elementary elements can prove the
inverse identities; the uniqueness clause often proves the same equality
more directly. Chapter~13 remains independent of tensor quotients, while
the final duality and shuffle comparison explain why the two constructions
agree. This order of proof will remain visible at each layer.

For a selective comparison of proof architecture, see
Dummit--Foote~\cite{dummit-foote}, Sections 10.4 and 11.4--11.5.
The displayed expansions and explicit factorization checks here retain
a slower pace and specialize throughout to vector spaces over fields.

\section{Quotient Vector Spaces and Universal Factorization}
\label{sec:quotient-vector-spaces}
In $E=F^2$, suppose horizontal displacements are to be ignored. Then
$(x,y)$ and $(x+t,y)$ should describe the same object, a horizontal line.
The second coordinate survives, whereas the first does not. This resembles
the equivalence-class constructions of integers and rationals in Chapter~2 (\cref{prop:integer-arithmetic,prop:rational-arithmetic}):
different representatives describe a single element. Here the differences
we ignore form a vector subspace.
\begin{definition}[Cosets]
For a subspace $R\subseteq E$, declare
\[
x\sim y\quad\text{ if }\quad x-y\in R.
\]
The class of $x$ is the coset $x+R=\{x+r:r\in R\}$. Write $E/R$ for
the set of these classes.
\end{definition}
The subspace laws give reflexivity, symmetry, and transitivity: differences
are zero, negated, or added. Moreover $x+R=y+R$ if and only if $x-y\in R$,
since adding a vector of $R$ leaves its coset unchanged.
\begin{theorem}[The quotient vector space]\label{thm:quotient-space}
The operations
\[
(x+R)+(y+R)=(x+y)+R,\qquad a(x+R)=ax+R
\]
make $E/R$ a vector space. The quotient map $\pi:E\to E/R$, $\pi(x)=x+R$,
is linear and surjective, with kernel $R$.
\end{theorem}
\begin{proof}
Changing representatives to $x+r,y+s$ changes their sum by $r+s\in R$
and a scalar multiple by $ar\in R$. Thus the operations are well defined.
The vector-space identities descend from $E$; zero is $R$ and the inverse
of $x+R$ is $-x+R$. The displayed operations make $\pi$ linear. Every class
has a representative, and $\pi(x)=R$ exactly when $x\in R$.
\end{proof}

\begin{example}[A legal rule and an illegal rule]
For $R=\{(t,0):t\in F\}\subseteq F^2$, the rule
$\mu((x,y)+R)=y$ is independent of representative. Its reverse is
$\eta(y)=(0,y)+R$. Both are linear by the quotient operations;
$\mu\eta(y)=y$, and $\eta\mu((x,y)+R)=(0,y)+R=(x,y)+R$ because their
difference is horizontal. Thus $F^2/R\cong F$ by these specific maps.
The rule $(x,y)+R\mapsto x$ is not a function on classes: $(0,0)$ and
$(1,0)$ represent the same class but would give different outputs.
\end{example}

\begin{theorem}[Quotient universal factorization]\label{thm:quotient-factorization}
For a linear $L:E\to U$, a linear map $\widetilde L:E/R\to U$ satisfying
$L=\widetilde L\circ \pi$ exists if and only if $R\subseteq\ker L$.
When it exists it is unique and is given by
\[
\widetilde L(x+R)=L(x).
\]
\end{theorem}
\[
\begin{tikzcd}[column sep=large,row sep=large]
E\arrow[r,"\pi"]\arrow[dr,"L"']&E/R\arrow[d,dashed,"\widetilde L"]\\
&U
\end{tikzcd}
\]
The diagram says both routes starting at $x\in E$ give $L(x)\in U$.
All three arrows here are linear; the dashed one is the map to be constructed.
\begin{proof}
If $L=\widetilde L\pi$ and $r\in R$, then $\pi(r)=0$ and
$L(r)=\widetilde L(0)=0$, proving necessity.
For sufficiency suppose $R\subseteq\ker L$. If $x+R=y+R$, then
$x-y\in R$, so $L(x)-L(y)=L(x-y)=0$. The displayed rule is therefore
independent of representative and defines $\widetilde L$.
For classes $x+R,y+R$ and scalars $a,b$,
\[
\widetilde L\bigl(a(x+R)+b(y+R)\bigr)=L(ax+by)
=a\widetilde L(x+R)+b\widetilde L(y+R).
\]
Thus it has the required type, a linear map on the quotient, and
$\widetilde L(\pi(x))=L(x)$. For uniqueness, any map $H$ with $H\pi=L$ has
$H(x+R)=H(\pi(x))=L(x)$ on every class, since $\pi$ is surjective. Hence
$H=\widetilde L$.
\end{proof}
The phrase ``$L$ descends to the quotient'' abbreviates this construction
under the relation-killing hypothesis.

\begin{example}[First isomorphism theorem]\label{thm:first-isomorphism}
For linear $T:V\to W$, quotient factorization gives
\[
V/\ker T\longrightarrow\operatorname{im}T,\qquad v+\ker T\longmapsto Tv.
\]
Its kernel is zero and it is onto its stated target, so it is an isomorphism.
Its inverse sends $Tv$ to $v+\ker T$: any two preimages differ by a kernel
vector. For $T(x,y,z)=(x-y,y-z)$ the kernel consists of $(t,t,t)$;
the class of $(a,0,-b)$ corresponds to $(a,b)$.
\end{example}

\begin{proposition}[Quotient bases and chosen complements]\label{prop:quotient-basis}
Let $E$ be finite dimensional and $R\subseteq E$ a subspace. Extend a
basis $r_1,\ldots,r_d$ of $R$ to $r_1,\ldots,r_d,s_1,\ldots,s_m$ of $E$.
Then $s_1+R,\ldots,s_m+R$ form a basis of $E/R$, and
\[
\dim(E/R)=\dim E-\dim R.
\]
For any chosen complement $E=R\oplus S$, restriction gives an isomorphism
$\pi|_S:S\to E/R$.
\end{proposition}
\begin{proof}
Expand $x\in E$ in the extended basis. Its $r$-part disappears in the
quotient, proving spanning of the surviving cosets. If
$\sum_j a_j(s_j+R)=0$, then $\sum_j a_js_j\in R$. Expanding that vector
in the $r_i$ and using independence of the full basis gives every $a_j=0$.
This proves independence and the dimension formula separately.
For a complement, $\pi|_S$ is linear, has kernel $S\cap R=0$, and is onto
since $\pi(r+s)=\pi(s)$. Its reverse sends $x+R$ to the $S$-component of $x$:
adding an element of $R$ does not change that component. Componentwise
linearity makes the reverse linear; the composites return $s$ and $x+R$.
\end{proof}
The quotient is not literally its chosen complement. Different complements
give different representative choices; the quotient map chooses none canonically.

\section{Bilinear Maps and Binary Tensor Products}
A bilinear rule accepts two independent vectors. Scalar multiplication
$(a,v)\mapsto av$, evaluation $(\lambda,v)\mapsto\lambda(v)$, and the
coordinate table $(x,y)\mapsto(x_iy_j)_{ij}$ are examples. The word
``independent'' describes how the arguments may be varied, not linear
independence of the vectors. Fixing either argument leaves a linear map
in the other; varying both together need not do so.
For example $b:F\times F\to F$, $b(x,y)=xy$, has
$b((1,0)+(0,1))=1$ and $b(1,0)+b(0,1)=0$. Thus the Cartesian product,
with its usual vector addition, is not the source that linearizes this rule.
We seek a vector space representing \emph{all} bilinear rules on the pair.
\begin{definition}[Bilinearity]
A map $b:V\times W\to U$ is bilinear if, for all vectors and scalars,
\begin{align*}
b(v+v',w)&=b(v,w)+b(v',w),\\
b(av,w)&=ab(v,w),\\
b(v,w+w')&=b(v,w)+b(v,w'),\\
b(v,aw)&=ab(v,w).
\end{align*}
These are two independent linearity requirements, one per slot.
\end{definition}
A tensor product retains the common slot identities of these examples,
without imposing special identities of one particular rule such as a dot product.
\begin{definition}[Tensor product]\label{def:tensor-universal}
A tensor product of $V,W$ consists of a vector space $T$ and a bilinear map
$\tau:V\times W\to T$ with the following property:
\begin{enumerate}
\item For every vector space $U$ over $F$,
\item and every bilinear map $b:V\times W\to U$,
\item there exists a unique linear map $\widetilde b:T\to U$ such that
$\widetilde b\circ\tau=b$.
\end{enumerate}
Write $T=V\otimes_F W$ and $v\otimes w=\tau(v,w)$; the latter is an
elementary tensor. We omit the subscript $F$ when the field is understood.
\end{definition}
\[
\begin{tikzcd}[column sep=large,row sep=large]
V\times W\arrow[r,"\tau"]\arrow[dr,"b"']&V\otimes W\arrow[d,dashed,"\widetilde b"]\\
&U
\end{tikzcd}
\]
``Every'' lets us choose $U$ and $b$ for the problem at hand. ``Exists''
gives an actual linear map on tensors, not just a formula on a pair.
``Unique'' says two linear maps with this same composite with $\tau$ must
agree. The horizontal and diagonal arrows are bilinear; the vertical one is
linear. The equation means both routes send $(v,w)$ to $b(v,w)$.

\subsection{Construction from formal symbols}

Let $E$ consist of finitely supported coefficient families indexed by pairs
$(v,w)\in V\times W$. Write such a family as $\sum c_{v,w}[v,w]$.
Here finite support means the set $\{(v,w):c_{v,w}\ne0\}$ is finite;
no closure or topology is involved. The pair $(v,w)$ is a label,
$[v,w]$ its basis vector, $c_{v,w}$ a scalar, and the full sum an element of $E$.
The symbol $[v,w]$ has coefficient $1$ at that pair and zero elsewhere.
These symbols form a basis by the definition of coefficient families;
addition and scaling act on their coefficients. In particular $[v+v',w]$
is initially a single symbol, not the sum $[v,w]+[v',w]$. Even $[0,w]$
is a nonzero basis symbol at this stage.

Let $R$ be the subspace spanned by all vectors of the following four types:
\begin{align*}
&[v+v',w]-[v,w]-[v',w],\\
&[av,w]-a[v,w],\\
&[v,w+w']-[v,w]-[v,w'],\\
&[v,aw]-a[v,w].
\end{align*}
Define $T=E/R$ and $\tau(v,w)=[v,w]+R$. Each displayed difference becomes
zero in the quotient, so $\tau$ satisfies the four bilinear identities.
No other types of relations have been imposed.
\begin{theorem}[Existence of the tensor product]\label{thm:tensor-existence}
The pair $(E/R,\tau)$ satisfies the tensor universal property.
\end{theorem}
\begin{proof}
Fix an arbitrary target $U$ and bilinear $b:V\times W\to U$.
\emph{First extend from the formal basis.} Define
\[
\widehat b:E\to U,\qquad
\widehat b\left(\sum c_{v,w}[v,w]\right)=\sum c_{v,w}b(v,w).
\]
Each sum is finite, and its coefficients are uniquely specified in $E$.
The formula is linear by distributing coefficients. At this stage
bilinearity of $b$ has not been used: any assignment of values to formal
basis symbols extends linearly in this way.

\emph{Then check the relations.} Bilinearity now gives all four checks:
\begin{align*}
\widehat b([v+v',w]-[v,w]-[v',w])&=b(v+v',w)-b(v,w)-b(v',w)=0,\\
\widehat b([av,w]-a[v,w])&=b(av,w)-ab(v,w)=0,\\
\widehat b([v,w+w']-[v,w]-[v,w'])&=b(v,w+w')-b(v,w)-b(v,w')=0,\\
\widehat b([v,aw]-a[v,w])&=b(v,aw)-ab(v,w)=0.
\end{align*}
Linearity makes $\widehat b$ vanish on their whole span:
$R\subseteq\ker\widehat b$. The quotient factorization theorem
\cref{thm:quotient-factorization} therefore gives the linear map
$\widetilde b:E/R\to U$, with
\[
\widetilde b(e+R)=\widehat b(e),\qquad
\widetilde b(v\otimes w)=b(v,w).
\]
Thus $\widetilde b\tau=b$. Any other linear $H$ with this property has
$H\pi([v,w])=b(v,w)=\widehat b([v,w])$ on the formal basis, hence
$H\pi=\widehat b$. Quotient uniqueness gives $H=\widetilde b$.
\end{proof}
The two stages solve different problems: basis extension gives a map before
identifications; killing the relations allows it to survive identifications.
\[
\begin{tikzcd}[column sep=large]
E\arrow[r,"\pi"]\arrow[dr,"\widehat b"']&E/R\arrow[d,dashed,"\widetilde b"]\\
&U
\end{tikzcd}
\]

\begin{center}
\fbox{\parbox{.88\linewidth}{\textbf{Working rule.}
To construct a map from a tensor product: define a concrete map on the
independent variables, verify bilinearity, invoke the universal property,
and obtain the induced linear map.

To prove two existing linear maps equal: first establish both maps, compare
their composites with the bilinear input map, then use uniqueness (or,
equivalently here, agreement on the spanning elementary tensors). A formula on
elementary tensors alone is not a definition of a linear map.}}
\end{center}
\begin{example}[Construct before evaluating]
On $F^2\times F^2$ let $b(v,w)=v_1w_1+2v_1w_2-v_2w_1$.
For fixed $w$ it is $(w_1+2w_2)v_1-w_1v_2$, linear in $v$; for fixed
$v$ it is $(v_1-v_2)w_1+2v_1w_2$, linear in $w$. Thus a unique linear
$L:F^2\otimes F^2\to F$ exists with these input values. It sends
$3e_1\otimes e_2+e_2\otimes e_1$ to $6-1=5$, interpreted in $F$.
By contrast the proposed rule $v\otimes w\mapsto v$ fails: $v\otimes0=0$
for every $v$, whereas the proposed outputs vary. Its independent-variable
map is not linear in $w$, so the working rule detects the problem in advance.
\end{example}

\begin{theorem}[Uniqueness up to canonical isomorphism]\label{thm:tensor-uniqueness}
If $(T,\tau)$ and $(T',\tau')$ are tensor products of $V,W$, there is a
unique linear isomorphism $A:T\to T'$ satisfying $A\tau=\tau'$.
\end{theorem}
\begin{proof}
\emph{Construct $A$.} Use the universal property of $T$ with target
$U=T'$ and input bilinear map $b=\tau'$. It gives a unique linear
$A:T\to T'$ such that $A\tau=\tau'$. Temporarily distinguishing the
spaces, this reads $A(v\otimes_T w)=v\otimes_{T'}w$.

\emph{Construct the reverse.} Use the universal property of $T'$ with
target $T$ and bilinear map $\tau$. This gives a linear $C:T'\to T$
with $C\tau'=\tau$, or $C(v\otimes_{T'}w)=v\otimes_Tw$.

\emph{First composite.} Both $CA$ and $\operatorname{id}_T$ are existing
linear maps $T\to T$. Compute
\[
(CA)\tau=C(A\tau)=C\tau'=\tau=\operatorname{id}_T\tau.
\]
The uniqueness clause for $T$, with target $T$ and input map $\tau$, says
only one linear map has this composite. Therefore $CA=\operatorname{id}_T$.

\emph{Second composite.} Both $AC$ and $\operatorname{id}_{T'}$ are linear
maps $T'\to T'$, and
\[
(AC)\tau'=A(C\tau')=A\tau=\tau'=\operatorname{id}_{T'}\tau'.
\]
Uniqueness for $T'$, now with target $T'$ and input map $\tau'$, gives
$AC=\operatorname{id}_{T'}$. Hence $A,C$ are mutual inverses. The same
uniqueness used to construct $A$ proves that no other isomorphism satisfying
the stated input condition exists.
\end{proof}
\Needspace{9\baselineskip}
The map $A$ is canonical: its behavior on every pair uniquely determines it,
without any basis choice. Two constructions need not be literally equal
as sets to be identified by this particular isomorphism.
\[
\begin{tikzcd}[column sep=large]
T\arrow[rr,shift left=1ex,"A"]&&T'\arrow[ll,shift left=1ex,"C"]\\
&V\times W\arrow[ul,"\tau"]\arrow[ur,"\tau'"']&
\end{tikzcd}
\]


\paragraph{Working facts: elementary tensor arithmetic.}
Bilinearity gives
\[
0\otimes w=v\otimes0=0,\qquad
(av)\otimes w=a(v\otimes w)=v\otimes(aw).
\]
Expanding $(0+0)\otimes w$ and subtracting $0\otimes w$ gives zero;
the other slot is identical. The scalar equalities are the two scalar
bilinearity identities.

\subsection{Spanning is not surjectivity of the input map}
Every quotient class of a finite sum of symbols is a finite sum of
elementary tensors. Therefore elementary tensors span. The universal map
$\tau:V\times W\to V\otimes W$ need not be surjective: being a finite
sum of elementary tensors is weaker than being one elementary tensor.
\begin{proposition}[Universality itself forces spanning]
For any tensor product $(T,\tau)$, $\operatorname{span}\tau(V\times W)=T$.
\end{proposition}
\begin{proof}
Let $S$ be that span. The quotient map $\pi:T\to T/S$ has $\pi\tau=0$.
The zero linear map $Z:T\to T/S$ also has $Z\tau=0$. Tensor uniqueness
with target $T/S$ and zero bilinear input forces $\pi=Z$. But $\pi$ is
surjective, so $T/S$ is the zero space. Hence every $t\in T$ belongs
to $S$, as claimed.
\end{proof}
Representations as sums need not be unique. For instance
\[
(v+v')\otimes w=v\otimes w+v'\otimes w,\qquad
(av)\otimes w=v\otimes(aw).
\]
The elementary factors and even the number of summands can change. A rule
on a chosen representation therefore requires justification.


\section{Finite Associativity}
\label{sec:tensor-associativity}
Both parenthesizations solve the same problem: linearize a trilinear
map on $U\times V\times W$. This is the reason for associativity.
\begin{theorem}[Associativity]
There is a canonical isomorphism
\[
A:(U\otimes V)\otimes W\longrightarrow U\otimes(V\otimes W),\qquad
A((u\otimes v)\otimes w)=u\otimes(v\otimes w).
\]
\end{theorem}
\begin{proof}
We first verify the common universal problem. Given trilinear
$b:U\times V\times W\to Z$, fixing $w$ gives a bilinear map and hence
a unique linear $L_w:U\otimes V\to Z$ with $L_w(u\otimes v)=b(u,v,w)$.
Uniqueness in binary universality implies
$L_{aw+cw'}=aL_w+cL_{w'}$: both sides factor the same bilinear input map.
Consequently $(t,w)\mapsto L_w(t)$ is bilinear on $(U\otimes V)\times W$
and factors uniquely through $(U\otimes V)\otimes W$.
Fixing $u$ instead proves the identical trilinear universal property for
$U\otimes(V\otimes W)$.

Apply the first property to the trilinear input
$(u,v,w)\mapsto u\otimes(v\otimes w)$ to construct $A$; apply the second
to $(u,v,w)\mapsto(u\otimes v)\otimes w$ to construct its reverse $C$.
The composites $CA$ and $AC$ factor their respective original tuple maps,
as do the identity maps. Uniqueness in the two trilinear properties gives
$CA=\operatorname{id}$ and $AC=\operatorname{id}$.
\end{proof}
The finite version below uses the same argument inductively; no basis or
dimension calculation is involved.

\section{Symmetry over a Field}

The two inputs have occupied ordered slots, but we may wish to exchange
their roles. Exchanging slots changes the target tensor product; it does
not assert equality of tensors within the original space. The universal
property turns this exchange into an isomorphism.

\begin{proposition}[Symmetry]
The rule $v\otimes w\mapsto w\otimes v$ determines a canonical isomorphism
$S:V\otimes W\to W\otimes V$.
\end{proposition}
\begin{proof}
The concrete map $(v,w)\mapsto w\otimes v$ is bilinear by the two
tensor-slot identities. The working rule constructs $S$; interchanging
the slots again constructs $R:W\otimes V\to V\otimes W$.
Now $RS(v\otimes w)=v\otimes w$ and $SR(w\otimes v)=w\otimes v$.
Each composite and the relevant identity factor the same bilinear input
map; uniqueness proves both inverse identities. No choice of representation has been made.
\end{proof}

In a space with basis $e_1,e_2$, symmetry interchanges $e_1\otimes e_2$ and $e_2\otimes e_1$.
The tensor-basis theorem below will prove that these are distinct.
The field scalars may be moved from either tensor slot to the other.
Swapping the inputs therefore still gives a bilinear map over the same
field. Symmetry changes their order; associativity only changes grouping.

\section{Tensor Products and Finite Direct Sums}
\label{sec:tensor-direct-sums}
For a finite family $(V_i)_{i\in I}$, the direct sum
$\bigoplus_{i\in I}V_i$ is the vector space of tuples $(v_i)$, with
componentwise addition and scalar multiplication. The injection
$\iota_i:V_i\to\bigoplus_i V_i$ places a vector in its $i$th component.
Every tuple is $\sum_i\iota_i(v_i)$. Given linear maps $L_i:V_i\to X$,
Chapter~10's direct-sum construction gives the unique linear map
$L:\bigoplus_iV_i\to X$ with $L\iota_i=L_i$:
\[
L((v_i))=\sum_i L_i(v_i).
\]
All index sets in the direct-sum formulas below are finite.

\begin{theorem}[Finite direct sums]\label{thm:tensor-finite-sums}
For finite index sets $I,J$ there are canonical isomorphisms
\begin{align*}
\Phi:(\bigoplus_{i\in I}V_i)\otimes W&\longrightarrow
\bigoplus_{i\in I}(V_i\otimes W),\\
\Theta:V\otimes(\bigoplus_{j\in J}W_j)&\longrightarrow
\bigoplus_{j\in J}(V\otimes W_j).
\end{align*}
Their elementary values are
$\Phi((v_i)\otimes w)=(v_i\otimes w)_i$ and
$\Theta(v\otimes(w_j))=(v\otimes w_j)_j$.
\end{theorem}
\begin{proof}
\emph{Construct the first forward map.} The proposed independent-variable
map is $b((v_i),w)=(v_i\otimes w)_i$. Its output is a tuple in the stated finite direct sum. Tensor
linearity in each component gives
\[
b(a\mathbf v+c\mathbf v',w)=ab(\mathbf v,w)+cb(\mathbf v',w),\qquad
b(\mathbf v,aw+cw')=ab(\mathbf v,w)+cb(\mathbf v,w').
\]
It is bilinear and hence induces $\Phi$.

\emph{Construct its reverse.} For each $i$, the map
$(v,w)\mapsto\iota_i(v)\otimes w$ is bilinear, because $\iota_i$ is
linear and the tensor map is bilinear. It induces
$j_i:V_i\otimes W\to(\bigoplus_iV_i)\otimes W$.
The direct-sum construction just proved gives the linear map
\[
\Psi:\bigoplus_i(V_i\otimes W)\longrightarrow(\bigoplus_iV_i)\otimes W,
\qquad \Psi((t_i))=\sum_i j_i(t_i).
\]
For an elementary tensor in the source of $\Phi$,
\[
\Psi\Phi((v_i)\otimes w)=\sum_i\iota_i(v_i)\otimes w
=\left(\sum_i\iota_i(v_i)\right)\otimes w=(v_i)\otimes w.
\]
For an elementary tensor injected into summand $i$ of the other space,
$\Phi\Psi$ returns that same injected tensor. Such vectors span the
direct sum, and elementary tensors span the first space. Both composites
are therefore identities.

\emph{Use symmetry for the other factor.} Compose the symmetry
$V\otimes(\bigoplus_jW_j)\cong(\bigoplus_jW_j)\otimes V$,
the proved first-factor distributivity, and symmetry in every resulting
summand. The composite sends $v\otimes(w_j)$ to $(v\otimes w_j)_j$,
which is exactly $\Theta$. Its inverse is the reverse composite.
\end{proof}

\begin{corollary}[Binary distributivity]\label{thm:tensor-distributivity}
In particular,
\begin{align*}
(U\oplus V)\otimes W&\cong(U\otimes W)\oplus(V\otimes W),\\
W\otimes(U\oplus V)&\cong(W\otimes U)\oplus(W\otimes V).
\end{align*}
\end{corollary}
\begin{proof}
Take a two-element index set in the two constructions above.
Their inverses add the maps induced by the two component injections.
\end{proof}
\begin{corollary}[Distributivity in both factors]\label{cor:finite-tensor-distributivity}
For finite index sets $I,J$,
\[
(\bigoplus_iV_i)\otimes(\bigoplus_jW_j)
\cong\bigoplus_{(i,j)\in I\times J}(V_i\otimes W_j).
\]
\end{corollary}
\begin{proof}
Apply the first-factor theorem with $W=\bigoplus_jW_j$, then apply
the second-factor theorem inside each summand. Reindex the finite tuple
of tuples by $(i,j)$. The composite sends $(v_i)\otimes(w_j)$ to
$(v_i\otimes w_j)_{i,j}$; all constituent maps are already isomorphisms.
\end{proof}
\section{Finite Multilinear Tensor Products}
\label{sec:finite-tensors}
Three different input spaces arise when a trilinear rule accepts three
different kinds of data. Repeated copies of one space are a special case.
We first state the input problem, then solve it by recursively using
the binary construction. A direct tuple-symbol quotient is an alternative realization, not a
second theory to develop.
\begin{definition}[Multilinearity]
For $n\geq1$, a map $b:V_1\times\cdots\times V_n\to U$ is multilinear
if it is linear in each slot with every other input fixed. In slot $i$,
\[
\begin{split}
b(v_1,\ldots,av_i+cv_i',\ldots,v_n)
={}&a b(v_1,\ldots,v_i,\ldots,v_n)\\
&+c b(v_1,\ldots,v_i',\ldots,v_n).
\end{split}
\]
Here $n$ is the number of factors and $i$ specifies the varying slot.
\end{definition}
\begin{definition}[Finite tensor-product universal property]
A tensor product of $V_1,\ldots,V_n$ consists of a vector space $T$
and a multilinear map $\tau:V_1\times\cdots\times V_n\to T$ such that
every multilinear map $b:V_1\times\cdots\times V_n\to U$ factors
uniquely as $b=\widetilde b\circ\tau$ with $\widetilde b:T\to U$ linear.
The value $\tau(v_1,\ldots,v_n)$ is an elementary tensor.
\end{definition}
\[
\begin{tikzcd}[column sep=large,row sep=large]
V_1\times\cdots\times V_n\arrow[r,"\tau"]\arrow[dr,"b"']&
T\arrow[d,dashed,"\widetilde b"]\\&U
\end{tikzcd}
\]
The horizontal and slant arrows are multilinear; the vertical arrow is
linear. We will build the upper-right space so that this factorization
always exists and is unique.

\Needspace{12\baselineskip}
\subsection{Recursive existence from the binary product}
Set
\[
P_1=V_1,\qquad P_n=P_{n-1}\otimes V_n\quad(n\geq2),
\]
and define the nested tuple maps by
\[
\tau_1(v_1)=v_1,\qquad
\tau_n(v_1,\ldots,v_n)=\tau_{n-1}(v_1,\ldots,v_{n-1})\otimes v_n.
\]
Tensor-slot linearity and induction make $\tau_n$ multilinear.
Its elementary values span $P_n$: elementary $t\otimes v_n$ span
the binary product, and expanding $t$ in nested elementary tensors
from $P_{n-1}$ expands each such tensor into values of $\tau_n$.
This argument uses no bases or dimension formula.

\begin{theorem}[Recursive construction and universality]\label{thm:finite-tensors}
For every finite $n\geq1$, $(P_n,\tau_n)$ represents $n$-linear maps.
\end{theorem}
\begin{proof}
Use $(n-1)$-linear universality, prove linear dependence
on the last variable, then apply binary universality.
For $n=1$, the given linear map itself is the unique factor through
$\tau_1=\operatorname{id}_{V_1}$. Assume the result at $n-1$ and take
an arbitrary target vector space $U$ and multilinear
$b:V_1\times\cdots\times V_n\to U$.

\emph{Fix the last input.} For $v_n\in V_n$, define
$b_{v_n}(v_1,\ldots,v_{n-1})=b(v_1,\ldots,v_n)$.
It is $(n-1)$-linear. The induction hypothesis constructs a linear map
\[
L_{v_n}:P_{n-1}\longrightarrow U,\qquad
L_{v_n}(\tau_{n-1}(v_1,\ldots,v_{n-1}))=b(v_1,\ldots,v_n).
\]
\emph{Let that input vary.} For $v,w\in V_n$ and $a,c\in F$, compare
the already existing linear maps $L_{av+cw}$ and $aL_v+cL_w$.
On each spanning elementary value $t=\tau_{n-1}(v_1,\ldots,v_{n-1})$,
linearity of the last slot of $b$ gives
\[
L_{av+cw}(t)=aL_v(t)+cL_w(t).
\]
Spanning proves equality on all of $P_{n-1}$. Thus
$v\mapsto L_v$ is a linear map $V_n\to\mathcal L(P_{n-1},U)$.

\emph{Apply binary universality.} Now define
\[
\beta:P_{n-1}\times V_n\longrightarrow U,\qquad \beta(t,v)=L_v(t).
\]
For fixed $v$, it is linear in $t$ by construction of $L_v$. For
fixed $t$, it is linear in $v$ by the map equality just proved.
Hence $\beta$ is bilinear and induces $\widetilde b:P_n\to U$.
Its value on $\tau_n(v_1,\ldots,v_n)$ is $b(v_1,\ldots,v_n)$.
Any two such factors agree on the spanning values of $\tau_n$, so
are equal. This proves existence and uniqueness at $n$.
\end{proof}

\begin{remark}[An alternative realization]\label{prop:finite-tensor-comparison}
One may alternatively start from tuple symbols modulo the slot-linearity
relations. The same relation check as in the binary construction gives
the finite multilinear universal property. Thus this realization is
canonically isomorphic to $P_n$: factor each object's tuple map through
the other and use uniqueness to identify both composites with identities.
There is no need to repeat the formal-symbol construction.
\end{remark}

\begin{corollary}[Every finite parenthesization]\label{thm:all-parenthesizations}
Every recursively binary parenthesization of the ordered factors
$V_1,\ldots,V_n$ represents the same multilinear problem as $P_n$.
Every canonical comparison preserves the elementary tuple values, and
any two canonical reassociation routes with the same endpoints agree.
\end{corollary}
\begin{proof}
Induct on the number of factors. At the outermost binary split, fix the
inputs on the right, factor the multilinear inputs on the left, and then
factor on the right. The uniqueness clause on the left ensures that the
first factor depends multilinearly on the right inputs, just as in the
proof of \cref{thm:finite-tensors}. This gives a bilinear map on the two
parenthesized blocks, and binary universality gives the required linear
factor on their tensor product. The same successive uniqueness clauses
prove its uniqueness. Hence every parenthesization represents the same
problem. Factoring their tuple maps constructs comparisons in both
directions. Uniqueness proves they are inverse and that all reassociation
routes with the same endpoints agree.
\end{proof}
For $n=3$, the left-associated tuple map sends $(u,v,w)$ to
$(u\otimes v)\otimes w$; the associativity map sends this to
$u\otimes(v\otimes w)$. Both represent the same trilinear input rule.
We may consequently write $V_1\otimes\cdots\otimes V_n$ without
parentheses, using these particular comparisons.

\begin{remark}[Finite-factor scope]
All tensor products in this appendix have finitely many factors, although
the individual vector spaces need not be finite dimensional. The finite
induction justifies finite regrouping. Infinitely many tensor factors
require a separate construction; the argument above neither defines
such a product nor justifies infinite regrouping. That specialized subject
is outside this appendix.
\end{remark}
\begin{definition}[Tensor powers]
Write $V^{\otimes n}$ for $n$ repeated factors. Set $V^{\otimes0}=F$
and identify the one-factor product with $V$ using $P_1=V$. For exterior powers below, $k$ will
denote the degree, and $V^{\otimes k}$ has $k$ factors.
\end{definition}

\section{Tensoring Linear Maps and Concrete Calculations}

If each input vector is transformed by a linear map, the combined tensor
data should transform as well. We must construct that transformation on
the whole tensor product, so its value does not depend on a decomposition
into elementary tensors.

\begin{proposition}[Induced binary tensor homomorphisms]\label{prop:binary-tensor-maps}
Linear maps $T:V\to V'$ and $S:W\to W'$ induce a unique linear map
\[
T\otimes S:V\otimes W\longrightarrow V'\otimes W',\qquad
(T\otimes S)(v\otimes w)=T(v)\otimes S(w).
\]
This construction preserves identities and composition, and takes pairs
of isomorphisms to isomorphisms.
\end{proposition}
\begin{proof}
The concrete map $(v,w)\mapsto T(v)\otimes S(w)$ is bilinear: linearity
of $T$ or $S$ followed by the corresponding tensor-slot identity verifies
each slot. The working rule constructs $T\otimes S$ with the stated type.
For further linear $T':V'\to V''$, $S':W'\to W''$, both composites
$(T'\otimes S')(T\otimes S)$ and $(T'T)\otimes(S'S)$ send $v\otimes w$
to $T'T(v)\otimes S'S(w)$, hence factor the same bilinear map and agree
by uniqueness. Identities act as
identities. If $T,S$ are invertible, the already justified construction
of $T^{-1}\otimes S^{-1}$ and the composition rule give both inverse identities.
\end{proof}

\subsection{Induced maps in all the factors}
For linear maps $T_i:V_i\to W_i$, consider directly
\[
b:V_1\times\cdots\times V_n\longrightarrow\bigotimes_iW_i,
\qquad b(v_1,\ldots,v_n)=T_1v_1\otimes\cdots\otimes T_nv_n.
\]
Linearity of $T_i$ and of the $i$th tensor slot proves multilinearity.
The finite universal property therefore constructs the unique linear map
\[
\bigotimes_iT_i:\bigotimes_iV_i\longrightarrow\bigotimes_iW_i
\]
with these values. For $S_i:W_i\to X_i$, the two maps below factor the
same multilinear rule $(v_i)\mapsto\bigotimes_i S_iT_iv_i$; uniqueness gives
\[
\bigotimes_i(S_iT_i)=(\bigotimes_iS_i)(\bigotimes_iT_i).
\]
The identity map factors the original tuple map, so tensoring identities
gives the identity. For invertible $T_i$, tensoring their inverses and
using composition gives both inverse laws. The recursive binary construction
has the same tuple values and is therefore compatible with this direct
construction. For repeated $T$ write $T^{\otimes n}$, and set
$T^{\otimes0}=\operatorname{id}_F$.

\subsection{The scalar unit}

The bilinear operation $(a,v)\mapsto av$ already has its values in $V$.
Linearizing it should recover $V$, with the scalar $1$ providing the
reverse map. This identifies a specific isomorphism, rather than merely
comparing dimensions.

\begin{proposition}[Scalar multiplication linearized]\label{prop:tensor-unit}
There is a canonical isomorphism
\[
\mu:F\otimes V\longrightarrow V,\qquad \mu(a\otimes v)=av.
\]
\end{proposition}
\begin{proof}
The vector-space axioms make $b:F\times V\to V$, $b(a,v)=av$,
bilinear. For example $b(a,cv+dw)=c\,b(a,v)+d\,b(a,w)$.
This bilinear map induces the linear $\mu$. Define the reverse directly:
\[
\eta:V\longrightarrow F\otimes V,\qquad \eta(v)=1\otimes v.
\]
The second tensor slot
is linear, so $\eta$ is linear. For every $v$, $\mu\eta(v)=1v=v$.
For an elementary tensor,
\[
\eta\mu(a\otimes v)=1\otimes(av)=a(1\otimes v)=a\otimes v.
\]
The two existing linear maps $\eta\mu$ and the identity agree on a
spanning set, hence agree everywhere. Thus both composites are identities.
\end{proof}
In particular $F\otimes_F F$ is one-dimensional although the formal
space has one independent symbol for every pair of field elements.

Composing the already constructed symmetry $V\otimes F\to F\otimes V$
with $\mu$ gives the right unit $v\otimes a\mapsto av$.
Its inverse is $v\mapsto v\otimes1$: composing in either order fixes
vectors or elementary tensors, respectively.

\subsection{Tensor bases and coordinates}

The construction has used no basis. With bases chosen, it should produce
a finite table of coefficients indexed by pairs of basis vectors. Expansion
shows that these pairs span; independent coordinate detectors must still
show that none of those coefficients is redundant.

\begin{theorem}[Tensor basis]\label{thm:tensor-basis}
If $(e_i)_{i=1}^m$ and $(f_j)_{j=1}^n$ are bases of $V,W$, the tensors
$e_i\otimes f_j$ form a basis of $V\otimes W$. Its dimension is $mn$.
\end{theorem}
\begin{proof}
\emph{Spanning.} Write $v=\sum_i x_ie_i$ and $w=\sum_j y_jf_j$.
Bilinearity expands $v\otimes w=\sum_{i,j}x_iy_j(e_i\otimes f_j)$.
Since elementary tensors span, so does the displayed finite list.

\emph{Independence.} The coordinate maps $e^r:V\to F$ and $f^s:W\to F$
are already linear by uniqueness of coordinates in the two given bases.
Define the scalar-valued input map $b_{rs}(v,w)=e^r(v)f^s(w)$.
For example
\[
b_{rs}(av+bv',w)=(ae^r(v)+be^r(v'))f^s(w)
=ab_{rs}(v,w)+bb_{rs}(v',w),
\]
and $b_{rs}(v,aw+bw')=ab_{rs}(v,w)+bb_{rs}(v,w')$ follows by expanding
the second coordinate functional. Thus it is bilinear. Universality gives
a linear $\ell_{rs}:V\otimes W\to F$, for which
\[
\ell_{rs}(e_i\otimes f_j)=\delta_{ri}\delta_{sj}.
\]
Applying this existing map to $\sum c_{ij}e_i\otimes f_j=0$ gives
$c_{rs}=0$. All coefficients vanish, proving independence.
\end{proof}
We did not define $\ell_{rs}$ as ``the coefficient extractor'': uniqueness
of tensor coefficients was precisely what the proof had to establish.
The bilinear detector constructed it without assuming the conclusion.
\begin{proposition}[Nonzero elementary tensors]
For finite-dimensional $V,W$, $v\otimes w=0$ if and only if $v=0$ or $w=0$.
\end{proposition}
\begin{proof}
A zero factor gives zero by elementary arithmetic. Otherwise extend $v$
and $w$ separately to bases. Their tensor is one of the tensor-basis
vectors, so is nonzero.
\end{proof}
\begin{example}[A table, and a tensor that is not elementary]
Choosing bases identifies tensors with rectangular coefficient tables.
A rectangular coefficient table alone determines no canonical map
$V\to W$; later, evaluation of a covector factor in $V^*\otimes W$
will define such a map.
The maps sending $\sum c_{ij}e_i\otimes f_j$ to $(c_{ij})$ and sending
$(c_{ij})$ back to that sum are linear and inverse: each composite restores
every coefficient. An elementary tensor has table $(x_iy_j)$. For example
$(e_1+2e_2)\otimes(f_1-f_2)$ in $F^2\otimes F^2$ has table
$\left(\begin{matrix}1&-1\\2&-2\end{matrix}\right)$.
The tensor $e_1\otimes f_1+e_2\otimes f_2$ is not elementary: an
elementary representation would require $x_1y_1=x_2y_2=1$ and
$x_1y_2=x_2y_1=0$. The first two equations make all four factors nonzero,
contradicting the last two. This proves that $\tau$ need not be surjective.
\end{example}

\begin{example}[Six coordinates]
In $F^2\otimes F^3$, order the basis by
$(e_1\otimes f_1,e_1\otimes f_2,e_1\otimes f_3,
e_2\otimes f_1,e_2\otimes f_2,e_2\otimes f_3)$.
Then $(e_1+2e_2)\otimes(f_1-f_2+3f_3)$ has coordinates
$(1,-1,3,2,-2,6)$. Integers are interpreted in $F$.
The vector-space dimension is six, although a single elementary tensor's
coordinates are constrained to be the products $x_i y_j$.
\end{example}

\begin{corollary}[Finite tensor basis]\label{cor:finite-tensor-basis}
Let $n\geq1$ and let
$\mathcal B_i=(e_1^{(i)},\ldots,e_{d_i}^{(i)})$ be a basis of $V_i$,
where $d_i=\dim V_i$. Then
\[
e_{j_1}^{(1)}\otimes\cdots\otimes e_{j_n}^{(n)},\qquad
1\leq j_i\leq d_i,
\]
form a basis of the finite tensor product, and its dimension is
$\prod_{i=1}^n d_i$.
\end{corollary}
\begin{proof}
For $P_1=V_1$, this is the given basis. If the nested basis is known
for $P_{n-1}$, apply the binary tensor-basis theorem
(\cref{thm:tensor-basis}) to $P_{n-1}\otimes V_n$.
It appends $e_{j_n}^{(n)}$ to each previous basis tensor and proves
both spanning and independence. Transport this basis through the canonical associativity maps of
\cref{thm:all-parenthesizations}. It becomes exactly the displayed list
in any parenthesized realization. Counting gives the product dimension.
If a factor is zero dimensional, the binary theorem gives the zero
space at that stage and thereafter, in agreement with the empty list.
\end{proof}
There is only one independent tensor-basis proof here, the binary one.
The recursive comparison was established without this corollary and
therefore legitimately transports the basis to any finite realization.

\begin{example}[A three-factor rule becomes a functional]\label{ex:three-factor-functional}
Consider three factor spaces with dimensions $2,2,3$:
\[
b:F^2\times F^2\times F^3\to F,\qquad
b(x,y,z)=x_1y_1z_1+x_2y_1z_2+x_1y_2z_3.
\]
Hold two inputs fixed. As a function of each remaining input the rule is
\begin{align*}
x&\longmapsto (y_1z_1+y_2z_3)x_1+(y_1z_2)x_2,\\
y&\longmapsto (x_1z_1+x_2z_2)y_1+(x_1z_3)y_2,\\
z&\longmapsto (x_1y_1)z_1+(x_2y_1)z_2+(x_1y_2)z_3.
\end{align*}
Each is a fixed linear combination of coordinate functions; substituting
$au+cv$ in its variable distributes into $a$ times the value at $u$ and
$c$ times the value at $v$. Thus all three linearity requirements hold.
With target $F$ and this input $b$, \cref{thm:finite-tensors} constructs
the unique linear functional
\[
\widetilde b:F^2\otimes F^2\otimes F^3\to F,
\qquad \widetilde b(x\otimes y\otimes z)=b(x,y,z).
\]
Write $e_i,f_j,g_k$ for the three standard bases. Its twelve basis values are
\[
\widetilde b(e_i\otimes f_j\otimes g_k)
=\delta_{i1}\delta_{j1}\delta_{k1}
+\delta_{i2}\delta_{j1}\delta_{k2}
+\delta_{i1}\delta_{j2}\delta_{k3}.
\]
Exactly the tensors indexed by $(1,1,1),(2,1,2),(1,2,3)$ have value $1$;
the other nine have value zero. For the finite sum
\[
t=e_1\otimes f_1\otimes g_1+e_2\otimes f_2\otimes g_2,
\qquad \widetilde b(t)=1+0=1.
\]
This $t$ is not elementary. In an elementary tensor the coefficients
$c_{ijk}=x_iy_jz_k$ satisfy $c_{111}c_{222}=c_{122}c_{211}$.
For $t$ the left side is $1$ and the right side is zero. Thus the new
functional evaluates an input that cannot be supplied as one triple to $b$;
its value on a sum is independent of the chosen representation.
\end{example}

\begin{example}[Four components of a tensor]\label{ex:four-tensor-components}
Take $A=F^2\oplus F$ and $B=F\oplus F^2$, with injections
$i:F^2\to A$, $j:F\to A$, $r:F\to B$, and $s:F^2\to B$.
Name the forward isomorphism in \cref{cor:finite-tensor-distributivity}
$D$ and its constructed inverse $E$. Thus $D$ has target
\[
(F^2\otimes F)\oplus(F^2\otimes F^2)
\oplus(F\otimes F)\oplus(F\otimes F^2),
\]
in that order. Let $e_1,e_2$ be the basis of the first $F^2$ and $f_1,f_2$
that of the second. On one elementary tensor it gives
\[
D((e_1+e_2,1)\otimes(1,f_1))
=((e_1+e_2)\otimes1,\ (e_1+e_2)\otimes f_1,\ 1\otimes1,\ 1\otimes f_1).
\]
The inverse $E$ is the already constructed sum of maps induced by the
four pairs of injections. For the component tuple
\[
z=(e_1\otimes1,\ e_2\otimes f_1+e_1\otimes f_2,\ 1\otimes1,\ 1\otimes f_2)
\]
its value is the nontrivial sum
\begin{align*}
E(z)={}&i(e_1)\otimes r(1)+i(e_2)\otimes s(f_1)
+i(e_1)\otimes s(f_2)\\
&+j(1)\otimes r(1)+j(1)\otimes s(f_2).
\end{align*}
Applying $D$ to the five terms gives, respectively,
\[
(e_1\otimes1,0,0,0),\ (0,e_2\otimes f_1,0,0),\ (0,e_1\otimes f_2,0,0),
\]
\[
(0,0,1\otimes1,0),\ (0,0,0,1\otimes f_2).
\]
Their sum is $z$, checking $DE$ on this input. For the first displayed
elementary tensor, $ED$ expands its four terms and recombines them as
$(i(e_1+e_2)+j(1))\otimes(r(1)+s(f_1))$, the original tensor.
Dimension counting then gives
$3\cdot3=2+4+1+2$. The count checks the construction but does not specify it.
\end{example}

\section{Tensor--Hom Correspondences}
\subsection{A map whose values are maps}
Recall from \cref{def:linear-map-space,prop:linear-map-space-structure} that
$\mathcal L(V,W)$ is the vector space of linear maps with pointwise addition
and scalar multiplication. A map into this space
returns a whole function, which can then be evaluated at a vector.
\begin{proposition}[Currying]\label{prop:tensor-hom}
There are inverse linear maps
\[
C:\mathcal L(V\otimes W,U)\to\mathcal L(V,\mathcal L(W,U)),\qquad
D:\mathcal L(V,\mathcal L(W,U))\to\mathcal L(V\otimes W,U).
\]
The forward rule is $C(L)(v)(w)=L(v\otimes w)$.
\end{proposition}
\begin{proof}
First check the types in the forward rule:
\[
\begin{aligned}
C(L)&\in\mathcal L(V,\mathcal L(W,U)),\\
C(L)(v)&\in\mathcal L(W,U),\\
C(L)(v)(w)&\in U.
\end{aligned}
\]
For fixed $v$, $w\mapsto L(v\otimes w)$ is linear by second-slot
linearity and linearity of $L$. Thus $C(L)(v)$ has the stated inner type.
For all $w$, first-slot linearity gives
$C(L)(av+bv')(w)=aC(L)(v)(w)+bC(L)(v')(w)$. Equality at every $w$
proves equality of maps, so $v\mapsto C(L)(v)$ is linear. Finally
$C(aL+bL')(v)(w)=aC(L)(v)(w)+bC(L')(v)(w)$; hence $C$ itself is linear.

For the inverse, take a linear $A:V\to\mathcal L(W,U)$ and define
$b_A(v,w)=A(v)(w)$. Since $A$ is linear,
$b_A(av+bv',w)=aA(v)(w)+bA(v')(w)$. Since $A(v)$ is a linear map,
$b_A(v,aw+bw')=aA(v)(w)+bA(v)(w')$. Thus $b_A$ is bilinear and induces
the linear $D(A):V\otimes W\to U$. Comparing existing maps on elementary
tensors shows $D(aA+bA')=aD(A)+bD(A')$, so $D$ is linear as well.

For every $v,w$,
$C(D(A))(v)(w)=D(A)(v\otimes w)=A(v)(w)$, proving $CD(A)=A$.
For every elementary tensor,
$D(C(L))(v\otimes w)=C(L)(v)(w)=L(v\otimes w)$.
Spanning proves $DC(L)=L$. These are both inverse identities.
\end{proof}

\subsection{A covector and a vector make a linear map}

A covector can measure an input and use the resulting scalar to rescale
a fixed output vector. Thus $v\mapsto\lambda(v)w$ is a linear map
with image in the line spanned by $w$. In finite dimension, basis
coordinates let every linear map be assembled from finitely many such pieces.

\begin{proposition}[Finite-dimensional tensor--map identification]
If $V$ is finite dimensional, there is a canonical isomorphism
\[
J:V^*\otimes W\longrightarrow\mathcal L(V,W),\qquad
J(\lambda\otimes w)(v)=\lambda(v)w.
\]
\end{proposition}
\begin{proof}
Begin with $B(\lambda,w)$ equal to the map $v\mapsto\lambda(v)w$.
This is a linear map $V\to W$ because $\lambda$ is linear. Evaluation at
each $v$ shows
$B(a\lambda+b\mu,w)=aB(\lambda,w)+bB(\mu,w)$ and
$B(\lambda,aw+bw')=aB(\lambda,w)+bB(\lambda,w')$ as equalities of maps.
So $B$ is bilinear with target $\mathcal L(V,W)$ and induces $J$.

Choose a basis $(e_i)$ of $V$ and its dual basis $(e^i)$ as in
\cref{def:dual-basis} (the same coordinate argument works over $F$).
Define the reverse
by the finite sum
\[
K(T)=\sum_i e^i\otimes T(e_i).
\]
Pointwise linearity in $T$ and tensor linearity show that $K$ is linear.
For $v=\sum_i e^i(v)e_i$,
\[
J(K(T))(v)=\sum_i e^i(v)T(e_i)=T(v),
\]
so $JK=\operatorname{id}$. On an elementary tensor the other composite is
\[
K(J(\lambda\otimes w))=\sum_i e^i\otimes\lambda(e_i)w
=\left(\sum_i\lambda(e_i)e^i\right)\otimes w=\lambda\otimes w.
\]
The last covector equality follows by evaluation on every $e_j$.
Spanning now proves $KJ=\operatorname{id}$. Although a basis was used to
write $K$, it is the unique inverse of the basis-free map $J$. Any other
basis therefore writes the same inverse, not another isomorphism.
\end{proof}
Finite dimension is used in the finite expansion of the identity on $V$.
With a target basis $(f_i)$, $J(e^j\otimes f_i)$ sends $e_j$ to $f_i$
and the other basis vectors to zero: its matrix unit is in row $i$, column $j$.
For example $e^1\otimes w_1+e^2\otimes w_2$ represents the map
$x_1e_1+x_2e_2\mapsto x_1w_1+x_2w_2$, the basis-image rule of Chapter~10.


\paragraph{Naturality of currying.}
For linear $R:V'\to V$, $S:W'\to W$, $T:U\to U'$ and
$L:V\otimes W\to U$, the already constructed maps satisfy
\[
C\bigl(T\circ L\circ(R\otimes S)\bigr)(v')(w')
=T\bigl(C(L)(R(v'))(S(w'))\bigr).
\]
Both sides equal $T(L(R(v')\otimes S(w')))$. This checks compatibility
with composition by evaluating at the two independent inputs; it also
shows why the displayed correspondence requires no bases.

\section{Exterior Powers}
\label{sec:exterior-quotients}
\subsection{Begin with two inputs}
Signed area is bilinear in its two edge vectors and vanishes when they
agree. Ordinary tensors retain $v\otimes v$; to model area-like rules
we must impose a further relation. The problem is to find a vector space
with an alternating bilinear input map such that every alternating
bilinear rule $V\times V\to U$ factors uniquely to a linear map from it.
The quotient will solve this problem. Define the subspace:
\[
R_2=\operatorname{span}\{v\otimes v:v\in V\}\subseteq V\otimes V.
\]
Then set $\Lambda^2V=(V\otimes V)/R_2$ and let $u\wedge v$ denote the
class of $u\otimes v$. The two-input wedge map is bilinear, as a composite
of a bilinear tensor map and a linear quotient map, and $v\wedge v=0$.
Expanding the repeated input $u+v$ gives
\[
0=(u+v)\wedge(u+v)=u\wedge v+v\wedge u,
\]
so $u\wedge v=-v\wedge u$. In characteristic two the minus sign is a
plus sign, but the repeated-input-zero relation still must be imposed.

\[
\boxed{(ae_1+be_2)\wedge(ce_1+de_2)=(ad-bc)e_1\wedge e_2.}
\]
For example $(e_1+2e_2)\wedge(3e_1+e_2)=-5e_1\wedge e_2$.
Over $\R$ the coefficient records signed area. This is the two-dimensional
calculation to keep in mind when we reach the determinant line.

The relation must use all vectors, not only basis vectors. In $F^2$,
killing just $e_1\otimes e_1,e_2\otimes e_2$ leaves the nonzero class
of $(e_1+e_2)\otimes(e_1+e_2)$, since its two off-diagonal coordinates
are $1$. The correct $R_2$ has basis
\[
e_1\otimes e_1,\quad e_2\otimes e_2,\quad
e_1\otimes e_2+e_2\otimes e_1.
\]
The third vector is obtained by subtracting the first two from the repeated
tensor of $e_1+e_2$. Conversely every repeated tensor has coordinates
$a^2,b^2,ab$ in this list. The list is independent by its tensor
coordinates. Thus $\dim R_2=3$ and its quotient has dimension $4-3=1$, even
in characteristic two.

The determinant is already alternating bilinear, so its linear tensor
factor kills $R_2$ and factors once more to a linear map $\Lambda^2F^2\to F$.
That map sends $e_1\wedge e_2$ to $1$, proving this wedge is nonzero.
The general version of this two-stage argument follows next.

\begin{definition}[Alternating multilinear map]
A multilinear $A:V^k\to U$ is alternating if
\[
v_i=v_j\text{ for some }i\ne j\quad\Longrightarrow\quad A(v_1,\ldots,v_k)=0.
\]
For $k=1$ the extra condition is empty. For $k\geq2$, expansion with
$u+v$ in two equal slots proves the interchange sign rule, just as above.
\end{definition}
For general $k\geq1$, we seek an alternating multilinear map from $V^k$
whose linear factors represent every alternating rule into every vector
space. The following quotient is the solution, verified immediately below.
\begin{definition}[Exterior power]
For $k\geq1$, let $R_k\subseteq V^{\otimes k}$ be the span of elementary
tensors having equal vectors in two distinct positions (so $R_1=\{0\}$).
The space $R_k$ is the smallest linear subspace containing all elementary
tensors with repeated inputs. Every alternating multilinear rule must kill
these tensors, so quotienting by $R_k$ imposes exactly the relations forced
by alternation. Define
\[
\Lambda^kV=V^{\otimes k}/R_k,\qquad
\iota_k(v_1,\ldots,v_k)=\pi(v_1\otimes\cdots\otimes v_k).
\]
Write this value $v_1\wedge\cdots\wedge v_k$, an elementary wedge.
We identify $\Lambda^1V$ with $V$, and set $\Lambda^0V=F$.
\end{definition}
Notice that $R_k$ is a span. A relation is not required to be one repeated
elementary tensor; it may be any finite linear combination of such tensors.

\begin{theorem}[Exterior universal property]\label{thm:exterior-universal}
The map $\iota_k$ is alternating multilinear. For every vector space $U$
and every alternating multilinear $A:V^k\to U$, there exists a unique
linear $\widetilde A:\Lambda^kV\to U$ with $\widetilde A\iota_k=A$.
\end{theorem}
\[
\begin{tikzcd}[column sep=large,row sep=large]
V^k\arrow[r,"\tau_k"]\arrow[dr,"A"']&V^{\otimes k}\arrow[r,"\pi"]\arrow[d,"L"]
&\Lambda^kV\arrow[dl,dashed,"\widetilde A"]\\
&U&
\end{tikzcd}
\]
\begin{proof}
The composite $\pi\tau_k$ is multilinear because $\tau_k$ is and $\pi$ is
linear. A repeated input puts $\tau_k(v_1,\ldots,v_k)$ in $R_k$, so its
quotient class is zero. Thus $\iota_k$ has the claimed type.

\emph{Tensor stage.} Given $A$, use its multilinearity and
\cref{thm:finite-tensors} to construct a unique linear
$L:V^{\otimes k}\to U$ with $L\tau_k=A$.

\emph{Relation check.} On a generator $v_1\otimes\cdots\otimes v_k$
of $R_k$, two inputs agree. Therefore
$L(v_1\otimes\cdots\otimes v_k)=A(v_1,\ldots,v_k)=0$ by alternation.
Linearity makes $L$ vanish on every finite combination of these generators,
so $R_k\subseteq\ker L$.

\emph{Quotient stage.} Apply \cref{thm:quotient-factorization} to obtain
$\widetilde A:\Lambda^kV\to U$, given by $\widetilde A(t+R_k)=L(t)$.
The relation check is precisely its representative-independence proof.
Now $\widetilde A\iota_k=\widetilde A \pi\tau_k=L\tau_k=A$.

\emph{Two uniqueness clauses.} If $H:\Lambda^kV\to U$ is linear with
$H\iota_k=A$, then $H\pi$ and $L$ are linear maps on $V^{\otimes k}$ with
the same composite $A$ with $\tau_k$. Tensor uniqueness gives $H\pi=L$.
Quotient uniqueness then gives $H=\widetilde A$. Existence used first
multilinearity and then alternation; uniqueness uses the two stages in the
same order.
\end{proof}
Elementary wedges span because tensor elements are finite sums of elementary
tensors and $\pi$ is surjective. If another space with an alternating input
map has the same universal property, construct $A$ from its input map and
$C$ from the original one, as in \cref{thm:tensor-uniqueness}. Both
$CA\iota_k=\iota_k$ and $AC\iota_k'=\iota_k'$ hold, and the corresponding
uniqueness clauses force both composites to be identities. Thus exterior
powers too are determined up to the unique input-preserving isomorphism.

\begin{center}
\fbox{\parbox{.88\linewidth}{\textbf{Working rule for exterior powers.}
Start with a concrete map on $V^k$, check multilinearity and repeated-input
vanishing, then invoke the exterior universal property. Elementary wedges
suggest a formula; this construction makes it well defined. Once maps
exist, equality on elementary wedges suffices, because they span.}}
\end{center}

\begin{proposition}[Elementary wedge calculus]
The elementary wedge is linear in each vector slot, is zero on repeated
inputs, changes sign on an interchange, and satisfies
\[
v_{\sigma(1)}\wedge\cdots\wedge v_{\sigma(k)}
=\operatorname{sgn}(\sigma)\,v_1\wedge\cdots\wedge v_k.
\]
\end{proposition}
\begin{proof}
Write $W(v_1,\ldots,v_k)=\iota_k(v_1,\ldots,v_k)$. Multilinearity gives,
with every unshown slot fixed,
\[
W(\ldots,au+bv,\ldots)=aW(\ldots,u,\ldots)+bW(\ldots,v,\ldots).
\]
Taking $a=b=0$ proves zero-slot vanishing. Alternation gives zero whenever
two inputs agree. In any two chosen slots it consequently gives
\[
\begin{aligned}
0&=W(\ldots,u+v,\ldots,u+v,\ldots)\\
 &=W(\ldots,u,\ldots,v,\ldots)+W(\ldots,v,\ldots,u,\ldots),
\end{aligned}
\]
because the equal-$u$ and equal-$v$ terms vanish. Moving one term to the
other side proves the interchange rule, with no division by $2$.
Decompose $\sigma$ into transpositions; each contributes a minus sign.
The parity theorem in Appendix~\ref{app:determinants} identifies their
product with $\operatorname{sgn}(\sigma)$, independently of the decomposition.
\end{proof}

\subsection{Expand, remove repeats, sort, and collect}
For $n=3,k=2$, the full expansion is
\begin{align*}
&(a_1e_1+a_2e_2+a_3e_3)\wedge(b_1e_1+b_2e_2+b_3e_3)\\
&=a_1b_1e_1\wedge e_1+a_1b_2e_1\wedge e_2+a_1b_3e_1\wedge e_3\\
&\quad+a_2b_1e_2\wedge e_1+a_2b_2e_2\wedge e_2+a_2b_3e_2\wedge e_3\\
&\quad+a_3b_1e_3\wedge e_1+a_3b_2e_3\wedge e_2+a_3b_3e_3\wedge e_3\\
&=(a_1b_2-a_2b_1)e_1\wedge e_2
+(a_1b_3-a_3b_1)e_1\wedge e_3
+(a_2b_3-a_3b_2)e_2\wedge e_3.
\end{align*}
The three diagonal terms vanish. Sorting $e_2\wedge e_1$,
$e_3\wedge e_1$, and $e_3\wedge e_2$ changes each sign. Collecting the
remaining pairs gives the three displayed coefficients. This proves
spanning by the three increasing wedges; independence is a separate task.
\begin{theorem}[Exterior basis]\label{thm:appendix-exterior-basis}
For a basis $(e_1,\ldots,e_n)$, the wedges
$e_I=e_{i_1}\wedge\cdots\wedge e_{i_k}$ with $i_1<\cdots<i_k$ form
a basis of $\Lambda^kV$. Thus $\dim\Lambda^kV=\binom nk$, and the space
is zero for $k>n$.
\end{theorem}
\begin{proof}
\emph{Spanning.} By \cref{cor:finite-tensor-basis}, the tensors
$e_{i_1}\otimes\cdots\otimes e_{i_k}$ span $V^{\otimes k}$.
The surjective quotient map $\pi:V^{\otimes k}\to\Lambda^kV$ sends them
to $e_{i_1}\wedge\cdots\wedge e_{i_k}$, so these wedges span the quotient.
Repeated indices vanish. For distinct indices, sort them into
$j_1<\cdots<j_k$, writing $i_s=j_{\sigma(s)}$; the interchange law gives
$e_{i_1}\wedge\cdots\wedge e_{i_k}
=\operatorname{sgn}(\sigma)e_{j_1}\wedge\cdots\wedge e_{j_k}$.
Thus only increasing wedges are needed. The preceding three-dimensional
expansion illustrates this general quotient-and-sorting mechanism.

\emph{Independence.} For an increasing tuple $I=(i_1,\ldots,i_k)$ define
the concrete map
\[
A_I(v_1,\ldots,v_k)=\det(e^{i_r}(v_s))_{r,s=1}^k.
\]
Changing $v_s$ linearly changes column $s$ linearly, so determinant
multilinearity gives linearity in that slot. Equal inputs give equal
columns and determinant zero. The determinant theory of Chapter~10 and
Appendix~\ref{app:determinants} therefore proves alternating multilinearity over $F$.
The exterior universal property gives a unique linear map
$\widetilde A_I:\Lambda^kV\to F$. For $J=(j_1<\cdots<j_k)$, its value is
\[
\widetilde A_I(e_J)=A_I(e_{j_1},\ldots,e_{j_k})
=\det(\delta_{i_rj_s})_{r,s=1}^k.
\]
If $I=J$ the matrix
is the identity, so $\widetilde A_I(e_I)=1$. If $I\ne J$, choose $j_s\in J\setminus I$, which exists since the two
different sets have the same size. The dual-basis rule
(\cref{def:dual-basis}) gives
\[
e^{i_r}(e_{j_s})=\delta_{i_rj_s}=0\quad(1\leq r\leq k).
\]
Thus column $s$ of the determinant is zero and $\widetilde A_I(e_J)=0$.
We have therefore proved $\widetilde A_I(e_J)=\delta_{IJ}$.
For a relation $\sum_Jc_Je_J=0$, fix $I$ and apply this constructed map:

\[
0=\widetilde A_I\Bigl(\sum_Jc_Je_J\Bigr)=\sum_Jc_J\widetilde A_I(e_J)=c_I.
\]
All coefficients vanish. There are $\binom nk$ increasing index sets;
for $k>n$ none remain and the spanning space is zero. In degree zero the
empty wedge is $1\in F$, giving dimension one.
\end{proof}
For instance $A_{13}(u,v)=u_1v_3-u_3v_1$ induces values $0,1,0$ on
$e_1\wedge e_2,e_1\wedge e_3,e_2\wedge e_3$. The induced linear map extracts a coefficient; its construction above
ensures that it respects every presentation of an exterior element.

\paragraph{Coordinates are minors.}
Write $v_j=\sum_{i=1}^n a_{ij}e_i$ for $1\leq j\leq k$. The linear maps just constructed give the full coordinate expansion
\[
v_1\wedge\cdots\wedge v_k
=\sum_{i_1<\cdots<i_k}\det(a_{i_rj})_{r,j=1}^k
\,e_{i_1}\wedge\cdots\wedge e_{i_k}.
\]
Indeed applying $\widetilde A_I$ to the left yields its displayed coefficient;
the exterior basis makes these coefficients unique. For two inputs these
are exactly the $2$ by $2$ minors obtained by expanding and collecting above.

\begin{proposition}[Nonvanishing of an elementary wedge]\label{prop:wedge-nonvanishing}
For finite-dimensional $V$, the wedge $v_1\wedge\cdots\wedge v_k$ is
nonzero exactly when its inputs are linearly independent.
\end{proposition}
\begin{proof}
For a dependent nonempty list, choose a relation with nonzero coefficient
at position $j$ and solve it as $v_j=\sum_{i\ne j}c_iv_i$. Then
\[
v_1\wedge\cdots\wedge v_k
=\sum_{i\ne j}c_i\,
v_1\wedge\cdots\wedge\underset{\text{slot }j}{v_i}\wedge\cdots\wedge v_k=0.
\]
Each summand repeats $v_i$ in slots $i,j$; for a one-vector dependent
list the vector is zero and linearity gives the same conclusion. For an independent list, extend it to a basis;
the wedge is an exterior-basis vector, hence nonzero. For $k=0$ the empty
list is independent and its wedge is $1$; for $k>\dim V$ no independent
list exists.
\end{proof}
Unlike $v\otimes v$ for $v\ne0$, $v\wedge v$ vanishes. The criterion
concerns elementary wedges, not arbitrary sums of them.

\subsection{Representations need not be unique}
Even an elementary wedge has many presentations:
$u\wedge(v+cu)=u\wedge v$, and for $c\ne0$,
$(cu)\wedge(c^{-1}v)=u\wedge v$. A general sum of elementary wedges
need not reduce to one. In $F^4$ put $\omega=e_1\wedge e_2+e_3\wedge e_4$.
For an elementary $u\wedge v$ its coefficients $p_{ij}=u_iv_j-u_jv_i$
satisfy $p_{12}p_{34}-p_{13}p_{24}+p_{14}p_{23}=0$. To verify this without
assuming any extra algebra, expand:
\begin{align*}
p_{12}p_{34}&=u_1u_3v_2v_4-u_1u_4v_2v_3-u_2u_3v_1v_4+u_2u_4v_1v_3,\\
-p_{13}p_{24}&=-u_1u_2v_3v_4+u_1u_4v_3v_2+u_3u_2v_1v_4-u_3u_4v_1v_2,\\
p_{14}p_{23}&=u_1u_2v_4v_3-u_1u_3v_4v_2-u_4u_2v_1v_3+u_4u_3v_1v_2.
\end{align*}
Commutativity pairs each term with its negative. For $\omega$, the only
nonzero increasing coefficients are $p_{12}=p_{34}=1$, making that
expression $1$. Thus $\omega$ is not elementary, in any characteristic.
Coordinates in a fixed wedge basis are unique; arbitrary elementary
decompositions are not. The product construction below must therefore
respect relations among different decompositions.

\subsection{Exterior maps, minors, and determinants}

A linear transformation sends a list of edge vectors to another list.
We want its action on their oriented area or volume data, including sums
of wedges with nonunique presentations. The exterior universal property turns an alternating rule on vector
tuples directly into a well-defined linear map on exterior elements.

\begin{proposition}[Induced exterior homomorphisms]\label{prop:exterior-maps}
Let $T:V\to W$ be linear. For $k\geq1$ there is a unique linear map
\[
\Lambda^kT:\Lambda^kV\longrightarrow\Lambda^kW,
\]
characterized by
\[
\Lambda^kT(v_1\wedge\cdots\wedge v_k)=Tv_1\wedge\cdots\wedge Tv_k.
\]
Set $\Lambda^0T=\operatorname{id}_F$. These induced maps commute with the
tensor quotient maps and preserve identities:
\[
\Lambda^k(\operatorname{id}_V)=\operatorname{id}_{\Lambda^kV}.
\]
For linear $S:W\to X$, they preserve composition:
\[
\Lambda^k(S\circ T)=(\Lambda^kS)\circ(\Lambda^kT).
\]
If $T$ is invertible, its induced map is invertible, with
\[
(\Lambda^kT)^{-1}=\Lambda^k(T^{-1}).
\]
\end{proposition}
\begin{proof}
We first construct a map on the direct product $V^k$, then use the
exterior universal property; no tensor representative needs to be chosen.
For $k\geq1$ put
\[
A_T:V^k\to\Lambda^kW,\qquad
A_T(v_1,\ldots,v_k)=Tv_1\wedge\cdots\wedge Tv_k.
\]
For a representative slot $j$, linearity of $T$ and that wedge slot gives
\[
\begin{aligned}
&A_T(v_1,\ldots,av_j+bv'_j,\ldots,v_k)\\
&=Tv_1\wedge\cdots\wedge T(av_j+bv'_j)\wedge\cdots\wedge Tv_k\\
&=Tv_1\wedge\cdots\wedge(aTv_j+bTv'_j)\wedge\cdots\wedge Tv_k\\
&=aA_T(v_1,\ldots,v_j,\ldots,v_k)
+bA_T(v_1,\ldots,v'_j,\ldots,v_k).
\end{aligned}
\]
If $v_i=v_j$, then $Tv_i=Tv_j$, so the output wedge is zero.
Thus $A_T$ is alternating multilinear. By \cref{thm:exterior-universal}
there is a unique linear $\Lambda^kT:\Lambda^kV\to\Lambda^kW$ with
\[
\Lambda^kT(v_1\wedge\cdots\wedge v_k)=A_T(v_1,\ldots,v_k).
\]
The construction establishes the map before we compare its values.
For $\xi=v_1\wedge\cdots\wedge v_k$,
\[
\Lambda^k(\operatorname{id}_V)(\xi)=v_1\wedge\cdots\wedge v_k=\xi,
\]
and
\[
\begin{aligned}
\Lambda^k(S\circ T)(\xi)
&=STv_1\wedge\cdots\wedge STv_k\\
&=(\Lambda^kS)(Tv_1\wedge\cdots\wedge Tv_k)\\
&=((\Lambda^kS)\circ(\Lambda^kT))(\xi).
\end{aligned}
\]
Each pair consists of linear maps factoring the same alternating input
map on $V^k$. Uniqueness in the exterior universal property therefore
proves both laws. If $T$ is invertible,
\[
(\Lambda^k(T^{-1}))(\Lambda^kT)=\Lambda^k(T^{-1}T)=\operatorname{id},
\]
and reversing the factors gives the identity on $\Lambda^kW$.
Hence $\Lambda^k(T^{-1})$ is the inverse. In degree zero all maps are
$\operatorname{id}_F$, so these statements hold directly.
\end{proof}

\paragraph{Compatibility with the quotient realization.}
We now compare the constructed map with tensors, rather than construct it
a second time. The tensor map $T^{\otimes k}$ sends a repeated-input
generator to a repeated-input tensor, since $v_i=v_j$ implies $Tv_i=Tv_j$.
For a finite linear combination of such generators $r_\ell$,
\[
T^{\otimes k}\Bigl(\sum_\ell c_\ell r_\ell\Bigr)
=\sum_\ell c_\ell T^{\otimes k}(r_\ell)\in R_k(W).
\]
Thus $T^{\otimes k}(R_k(V))\subseteq R_k(W)$. More explicitly, on every
elementary tensor $t=v_1\otimes\cdots\otimes v_k$ both routes give
\[
(\Lambda^kT)\pi_V(t)=Tv_1\wedge\cdots\wedge Tv_k
=\pi_WT^{\otimes k}(t).
\]
Both routes are linear and elementary tensors span, so the equality holds
on all tensors. This proves the compatibility asserted in the proposition:
\[
\begin{tikzcd}[column sep=large,row sep=large]
V^{\otimes k}\arrow[r,"T^{\otimes k}"]\arrow[d,"\pi_V"']&
W^{\otimes k}\arrow[d,"\pi_W"]\\
\Lambda^kV\arrow[r,"\Lambda^kT"']&\Lambda^kW.
\end{tikzcd}
\]
The square says that the direct universal-property construction agrees
with the concrete quotient realization.

\begin{proposition}[Minors of an exterior homomorphism]\label{prop:exterior-minors}
Let $V,W$ have bases $(e_j)_{j=1}^n,(f_i)_{i=1}^m$, and let the
possibly rectangular $m$-by-$n$ matrix $A$ represent $T:V\to W$.
For increasing $k$-tuples $J$ in the source and $I$ in the target,
\[
\Lambda^kT(e_J)=\sum_I\det(A_{I,J})f_I.
\]
\end{proposition}
\begin{proof}
We extract each coefficient by the linear map $\widetilde A_I$ constructed in
the exterior basis proof, now using the target dual basis $(f^i)$.
The column convention for $A$ says
\[
Te_{j_s}=\sum_r a_{rj_s}f_r,\qquad
f^{i_r}(Te_{j_s})=a_{i_rj_s}.
\]
Consequently the evaluation matrix is exactly the indicated minor:
\[
\bigl(f^{i_r}(Te_{j_s})\bigr)_{r,s}
=(a_{i_rj_s})_{r,s}=A_{I,J}.
\]
Writing $\Lambda^kT(e_J)=\sum_I c_I f_I$, we obtain
\[
c_I=\widetilde A_I(\Lambda^kT(e_J))
=\det\bigl(f^{i_r}(Te_{j_s})\bigr)_{r,s}=\det A_{I,J}.
\]
Substituting these coefficients proves the formula. For $k=0$ the empty
minor is $1$; above either dimension the relevant source or target is zero.
\end{proof}
For $k=1$ this is the usual basis-image formula $T(e_j)=\sum_i a_{ij}f_i$.
For example the columns $(1,0,1)$ and $(0,1,1)$ define $T:F^2\to F^3$.
Its three two-row minors are $1,1,-1$, so
$\Lambda^2T(e_1\wedge e_2)=f_1\wedge f_2+f_1\wedge f_3-f_2\wedge f_3$.

\begin{example}[Basis images still produce the matrix]
For $T(e_1)=e_1+e_2$, $T(e_2)=e_2$, $T(e_3)=2e_3$, the images of
$e_1\wedge e_2,e_1\wedge e_3,e_2\wedge e_3$ are respectively
$e_1\wedge e_2$, $2e_1\wedge e_3+2e_2\wedge e_3$, and
$2e_2\wedge e_3$. The matrix is
\[
[\Lambda^2T]=\begin{pmatrix}1&0&0\\0&2&0\\0&2&2\end{pmatrix}.
\]
This is the Chapter~10 basis-image procedure applied to an exterior basis.
\end{example}

The coefficient $ad-bc$ from our first two-input calculation now has an
intrinsic interpretation: the top exterior power is a line, and an
endomorphism acts on that line by one scalar.

\begin{theorem}[Fourth characterization of the determinant]
\label{thm:top-exterior-determinant}
\emph{Top exterior-power action.} If $\dim V=n$ and $T:V\to V$ is linear, then $\Lambda^nV$ is
one-dimensional and
\[
\Lambda^nT=(\det T)\operatorname{id}_{\Lambda^nV}.
\]
Equivalently, for every $v_1,\ldots,v_n\in V$,
\[
Tv_1\wedge\cdots\wedge Tv_n=(\det T)(v_1\wedge\cdots\wedge v_n).
\]
\end{theorem}
\begin{proof}
The top exterior power is a line. We identify its scalar action by applying
the normalized determinant form, whose value on the chosen generator is one.
Let $\mathcal B=(e_1,\ldots,e_n)$ and let
$\det_{\mathcal B}:V^n\to F$ be the already established alternating
$n$-linear map with $\det_{\mathcal B}(e_1,\ldots,e_n)=1$.
For $n\geq1$, the exterior universal property gives the unique linear map
\[
\widetilde{\det}_{\mathcal B}:\Lambda^nV\to F,\qquad
\widetilde{\det}_{\mathcal B}(v_1\wedge\cdots\wedge v_n)
=\det_{\mathcal B}(v_1,\ldots,v_n).
\]
By \cref{thm:appendix-exterior-basis}, $g=e_1\wedge\cdots\wedge e_n$
is nonzero and $\Lambda^nV=Fg$. Hence the endomorphism $\Lambda^nT$
satisfies $\Lambda^nT(g)=\lambda g$ for a unique $\lambda\in F$.
Now every step of the scalar identification is visible:
\[
\begin{aligned}
\lambda
&=\lambda\,\widetilde{\det}_{\mathcal B}(g)\\
&=\widetilde{\det}_{\mathcal B}(\lambda g)\\
&=\widetilde{\det}_{\mathcal B}(\Lambda^nT(g))\\
&=\widetilde{\det}_{\mathcal B}(Te_1\wedge\cdots\wedge Te_n)\\
&=\det_{\mathcal B}(Te_1,\ldots,Te_n)\\
&=\det[T]_{\mathcal B}=\det T.
\end{aligned}
\]
The first equality uses normalization, the second linearity, the next two
the defining action, and the fifth the factorization property. The last
line is the matrix expression and basis independence of the determinant
already established earlier. Both linear maps $\Lambda^nT$ and
$(\det T)\operatorname{id}$ agree on the basis element $g$ of this
one-dimensional space, so they agree on every $cg$. Only now apply this
map identity to arbitrary vectors:
\[
\begin{aligned}
Tv_1\wedge\cdots\wedge Tv_n
&=\Lambda^nT(v_1\wedge\cdots\wedge v_n)\\
&=(\det T)(v_1\wedge\cdots\wedge v_n).
\end{aligned}
\]
Replacing the generator by $g'=cg$, $c\ne0$, does not change the scalar:
\[
\Lambda^nT(g')=c\Lambda^nT(g)=c(\det T)g=(\det T)g'.
\]
Any generator from another ordered basis differs in this way, because
both are nonzero vectors in the same line. As a geometric consistency
check, if $T$ is singular its basis images are dependent, and
\[
Te_1\wedge\cdots\wedge Te_n=0
\]
by \cref{prop:wedge-nonvanishing}, matching $\det T=0$.
For $n=0$, $\Lambda^0V=F$, $\Lambda^0T=\operatorname{id}_F$, and
$\det T=1$, so the statement remains valid.
\end{proof}
For a coordinate confirmation, the minors formula has only one increasing
$n$-tuple in top degree. The matrix of $\Lambda^nT$ is therefore the
$1$-by-$1$ matrix $(\det[T]_{\mathcal B})$.

There are now four equivalent descriptions of the determinant:
normalized alternating multilinearity determines it; the Leibniz formula
constructs it; recursive Laplace expansion, starting with the empty
determinant $1$, computes that formula; and top exterior action records
the same scalar. These are established equivalences, not new axioms for
arbitrary computational algorithms. The exterior basis proof already uses
determinants, so this is a retrospective characterization.
In particular, the earlier multiplicativity theorem can now be understood
retrospectively through composition:
\[
\begin{aligned}
\Lambda^n(S\circ T)&=\det(S\circ T)\operatorname{id},\\
(\Lambda^nS)(\Lambda^nT)
&=(\det S)\operatorname{id}\circ(\det T)\operatorname{id}\\
&=(\det S)(\det T)\operatorname{id}.
\end{aligned}
\]
Functoriality gives equality of these two maps. Evaluating at the nonzero
generator $g$ forces equality of their scalar coefficients, hence
$\det(S\circ T)=\det S\det T$. This explains the earlier result; it
cannot serve as its foundational proof because the exterior basis argument
already used determinant theory.
For $T:V\to W$ between distinct $n$-dimensional spaces, the top map has
different source and target lines; there is no scalar without separate
basis or volume normalizations.

\subsection{Wedge products between exterior powers}
\label{sec:exterior-multiplication}
The intended elementary rule is
\[
(v_1\wedge\cdots\wedge v_p)\wedge(w_1\wedge\cdots\wedge w_q)
=v_1\wedge\cdots\wedge v_p\wedge w_1\wedge\cdots\wedge w_q.
\]
This is not yet an operation on general exterior elements: their expressions
as sums of elementary wedges are nonunique.

We want to concatenate two wedge lists, even when neither element is
given by a unique list. For $p,q\geq1$ the intended output lies in
$X=\Lambda^{p+q}V$. We will first construct a map from one exterior
power, then a map with values in a space of linear maps, and only then
evaluate to obtain a bilinear product.
\begin{lemma}[Equality on spanning tuples]\label{lem:multilinear-spanning}
If two existing multilinear maps $A,B:V_1\times\cdots\times V_r\to U$
agree on tuples from spanning sets $S_1,\ldots,S_r$, they agree everywhere.
\end{lemma}
\begin{proof}
Write $v_i=\sum_j a_{ij}s_{ij}$ with $s_{ij}\in S_i$. Expanding each
slot gives for either map the finite sum over index tuples with coefficient
$\prod_i a_{ij_i}$ and value at $(s_{1j_1},\ldots,s_{rj_r})$.
Those values agree by hypothesis, so the sums agree. The lemma compares
maps that already exist; it cannot prove that a proposed rule is a map.
\end{proof}
\begin{theorem}[Well-defined exterior product]
There is a unique bilinear map
$m_{p,q}:\Lambda^pV\times\Lambda^qV\to\Lambda^{p+q}V$ whose value on
two elementary wedges is their concatenated wedge.
\end{theorem}
\begin{proof}
\emph{First factor the $u$-variables.} Fix the tuple
$\mathbf v=(v_1,\ldots,v_q)$ and define
\[
b_{\mathbf v}(u_1,\ldots,u_p)
=u_1\wedge\cdots\wedge u_p\wedge v_1\wedge\cdots\wedge v_q\in X.
\]
The universal $(p+q)$-wedge map is linear in each $u$ slot, and equal
$u$ entries make its value zero. Thus $b_{\mathbf v}$ is alternating
$p$-linear. The exterior universal property constructs a linear
$L_{\mathbf v}:\Lambda^pV\to X$ with this value on elementary wedges.

The map
\[
B:V^q\to\mathcal L(\Lambda^pV,X),\qquad B(\mathbf v)=L_{\mathbf v},
\]
is alternating multilinear. Each value has the stated type by the preceding construction. Replace one
slot $v_j$ by $av_j+bv_j'$. There are now two existing linear maps
$L_{(\ldots,av_j+bv_j',\ldots)}$ and
$aL_{(\ldots,v_j,\ldots)}+bL_{(\ldots,v_j',\ldots)}$.
On an elementary $u_1\wedge\cdots\wedge u_p$, they agree by the
linearity of slot $p+j$ in the long wedge. Since elementary $p$-wedges
span, the two maps agree on every $\alpha\in\Lambda^pV$. This proves
linearity in slot $j$ of $B$, not just linearity of an isolated scalar value.
If two $v$ entries agree, $L_{\mathbf v}$ sends every elementary $p$-wedge
to zero by alternation of the long wedge. It is then the zero linear map
by spanning. Thus $B$ is alternating multilinear with values in
$\mathcal L(\Lambda^pV,X)$.

\emph{Factor the $v$-variables.} Apply the exterior universal property
again to obtain a linear map
\[
H:\Lambda^qV\to\mathcal L(\Lambda^pV,X),\qquad
H(v_1\wedge\cdots\wedge v_q)=L_{\mathbf v}.
\]
Define $m_{p,q}(\alpha,\beta)=H(\beta)(\alpha)$.
For fixed $\beta$, $H(\beta)$ is linear, giving linearity in $\alpha$.
For fixed $\alpha$, linearity of $H$ gives
\[
H(a\beta+b\beta')(\alpha)
=(aH(\beta)+bH(\beta'))(\alpha)
=aH(\beta)(\alpha)+bH(\beta')(\alpha).
\]
Thus $m_{p,q}$ is bilinear and has the prescribed concatenation values.
The construction does not require a unique elementary decomposition of a
general wedge.
For uniqueness, any other bilinear product with those values agrees on
pairs from the spanning elementary wedges; \cref{lem:multilinear-spanning}
gives equality everywhere.
\end{proof}
Write
$\alpha\wedge\beta=m_{p,q}(\alpha,\beta)$. When a degree is zero,
define the product by scalar multiplication in $F=\Lambda^0V$.
The construction can be tracked by the types
\[
V^q\longrightarrow\mathcal L(\Lambda^pV,\Lambda^{p+q}V)
\longleftarrow\Lambda^qV,
\]
where the left arrow is alternating multilinear and the right arrow is
the linear factor $H$, evaluated at $\alpha$ only at the final stage.

\begin{proposition}[Associativity and graded commutativity]
For homogeneous elements of degrees $p,q,r$,
\[
(\alpha\wedge\beta)\wedge\gamma=\alpha\wedge(\beta\wedge\gamma),
\qquad \alpha\wedge\beta=(-1)^{pq}\beta\wedge\alpha.
\]
\end{proposition}
\begin{proof}
Both expressions in associativity are existing trilinear maps of
$\alpha,\beta,\gamma$, since the products just constructed are bilinear.
On elementary wedges both give the same concatenation of the three lists.
The spanning-tuple lemma proves equality everywhere. For commutativity,
moving a block of $p$ vectors past a block of $q$ requires $pq$
interchanges, giving sign $(-1)^{pq}$ on elementary pairs. Both sides are
bilinear, so the same lemma extends this identity to arbitrary homogeneous
elements. A zero degree is governed by scalar multiplication, for which
associativity and the sign $1$ hold directly.
\end{proof}

\section{The Exterior Algebra and Alternating Forms}

Exterior multiplication has been constructed between separate degrees.
For a finite-dimensional space, only finitely many degrees are nonzero,
so we can collect them in one finite direct sum and multiply by distribution.
Afterward we relate these vector wedges to alternating scalar measurements.
\label{sec:appendix-exterior-algebra}
\begin{definition}[Exterior algebra]\label{def:appendix-exterior-algebra}
For a finite-dimensional vector space $V$ with $\dim V=n$, its
\textbf{exterior algebra} is the finite direct sum
\[
\Lambda V=\bigoplus_{k=0}^n\Lambda^kV,
\]
with multiplication obtained by distributing the constructed wedge products
between homogeneous components. The scalar $1\in\Lambda^0V=F$ is its unit.
\end{definition}
The exterior basis theorem makes degrees above $n$ zero, so this finite
sum contains every nonzero homogeneous component.
An element $\omega=\omega_0+\cdots+\omega_n$ has a unique component in
each degree, by the direct-sum construction. It is homogeneous of degree
$k$ if only that component can be nonzero. Extend the products above by
\[
\left(\sum_p\alpha_p\right)\wedge\left(\sum_q\beta_q\right)
=\sum_{p,q}\alpha_p\wedge\beta_q;
\]
degrees above $n$ are zero. The degree-$k$ component of the product is
$\sum_{p+q=k}\alpha_p\wedge\beta_q$. This finite sum is well defined
because the homogeneous components are unique. Associativity on homogeneous
elements and finite distribution prove associativity for arbitrary elements.
The scalar $1$ is an identity. ``Algebra'' here means a vector space with
this bilinear associative multiplication; no additional algebra course is
needed to use the construction.
\begin{example}[The complete exterior algebra of $F^3$]
For the standard basis, the degree table is
\[
\begin{array}{c|l|c}
\text{degree}&\text{basis}&\text{dimension}\\\hline
0&1&1\\
1&e_1,e_2,e_3&3\\
2&e_1\wedge e_2,e_1\wedge e_3,e_2\wedge e_3&3\\
3&e_1\wedge e_2\wedge e_3&1
\end{array}
\]
Thus the total dimension is $8$. Products are calculated by distribution:
\begin{align*}
(1+e_1)\wedge(1+e_2)&=1+e_1+e_2+e_1\wedge e_2,\\
(e_1\wedge e_2)\wedge e_3&=e_1\wedge e_2\wedge e_3,\\
e_2\wedge(e_1\wedge e_3)&=-e_1\wedge e_2\wedge e_3,\\
(e_1\wedge e_2)\wedge(e_2\wedge e_3)&=0.
\end{align*}
The first result has mixed degrees; the next two are homogeneous of degree
$3$. Graded commutativity applies directly to homogeneous factors.
Zero belongs to every homogeneous subspace and has no unique degree.
\end{example}

\begin{example}[Mixed degrees in dimension two]
For scalars $a,b$, vectors $u,v$, and degree-two elements $\omega_2,\eta_2$,
\[
(a+u+\omega_2)\wedge(b+v+\eta_2)
=ab+(av+bu)+(a\eta_2+b\omega_2+u\wedge v).
\]
The three groups have degrees $0,1,2$. Every omitted product has degree
above $2$ and vanishes. A mixed sum has no single homogeneous degree to
insert into the graded commutativity sign.
\end{example}
\begin{proposition}[The assembled exterior homomorphism]
Let $T:V\to W$ be linear, with $V,W$ finite dimensional. The direct sum of the maps $\Lambda^kT$
defines $\Lambda T:\Lambda V\to\Lambda W$, preserving degrees, the unit,
and wedge multiplication.
\end{proposition}
\begin{proof}
Each component map is linear and above the target dimension is zero;
thus the finite component rule defines a linear map. Degree zero fixes $1$.
For elementary wedges, applying $T$ before or after concatenation gives
the same list of images. The two existing bilinear maps agree on spanning
pairs, so agree on all homogeneous pairs. Finite distribution proves the
identity on mixed sums. This assembled product-preserving map is a graded
algebra homomorphism; each fixed-degree exterior homomorphism is a linear map.
\end{proof}

\begin{example}[Compute with basis wedges]
Write $e_I$ with its indices increasing. If $I,J$ overlap, $e_I\wedge e_J=0$;
otherwise sort their concatenation and attach its sign. Thus
\begin{align*}
&(2e_1+e_3)\wedge(e_1\wedge e_2+3e_2\wedge e_4)\\
&=6e_1\wedge e_2\wedge e_4+e_3\wedge e_1\wedge e_2
+3e_3\wedge e_2\wedge e_4\\
&=e_1\wedge e_2\wedge e_3+6e_1\wedge e_2\wedge e_4
-3e_2\wedge e_3\wedge e_4.
\end{align*}
The middle term needs two swaps, the last one swap. Over $\R$, the
bivector $\omega=e_1\wedge e_2+e_3\wedge e_4$ has
$\omega\wedge\omega=2e_1\wedge e_2\wedge e_3\wedge e_4\ne0$:
the two cross terms add because $(-1)^{2\cdot2}=1$.
In characteristic two this square is zero; the earlier coefficient
obstruction, not this square test, proves non-elementarity over every field.
\end{example}
\Needspace{10\baselineskip}
\paragraph{From alternating rules to linear maps into the field.}
\subsection{First identification: factor an alternating map}
Write $\operatorname{Alt}^k(V;F)$ for the vector space of alternating
$k$-linear scalar-valued maps, with pointwise operations; set $\operatorname{Alt}^0(V;F)=F$.
\begin{proposition}
There is a canonical isomorphism
$\operatorname{Alt}^k(V;F)\cong(\Lambda^kV)^*$.
\end{proposition}
\begin{proof}
For $A\in\operatorname{Alt}^k(V;F)$ use \cref{thm:exterior-universal} to
construct its unique linear factor $F(A)=\widetilde A:\Lambda^kV\to F$.
For the reverse define $G(\ell)(v_1,\ldots,v_k)=\ell(v_1\wedge\cdots\wedge v_k)$.
Linearity of $\ell$ preserves each slot's linearity and the zero value on
repeated inputs, so $G(\ell)$ is indeed in $\operatorname{Alt}^k(V;F)$.
Pointwise evaluation shows $G$ is linear. To check $F$, let
$\xi=v_1\wedge\cdots\wedge v_k$. Then
\[
F(aA+bB)(\xi)=(aA+bB)(v_1,\ldots,v_k)
=aF(A)(\xi)+bF(B)(\xi).
\]
The two sides are existing linear maps of $\xi$ agreeing on elementary
wedges; expansion of a finite sum of those wedges proves equality everywhere.
This proves linearity of $F$.
For each tuple, $G(F(A))(v_1,\ldots,v_k)=A(v_1,\ldots,v_k)$, so $GF$ is
the identity. On each elementary wedge,
$F(G(\ell))(v_1\wedge\cdots\wedge v_k)=\ell(v_1\wedge\cdots\wedge v_k)$.
Both are linear maps into $F$; spanning proves $FG$ is the identity.
In degree zero both spaces are $F$, with the identity identification.
\end{proof}

\subsection{Second identification: pair covector wedges with vector wedges}

The previous identification used a linear map from vector wedges to $F$.
Chapter~13 instead writes alternating forms as wedges of covectors.
To connect these two objects, evaluate a list of covectors on a list of
vectors and take the determinant. We must factor in each list separately;
finite-dimensional dual bases will then give the inverse.

\begin{theorem}[Exterior duality]\label{thm:exterior-duality}
For finite-dimensional $V$, there is a canonical isomorphism
$P:\Lambda^k(V^*)\to(\Lambda^kV)^*$ characterized by
\[
P(\lambda_1\wedge\cdots\wedge\lambda_k)
(v_1\wedge\cdots\wedge v_k)=\det(\lambda_i(v_j)).
\]
\end{theorem}
\begin{proof}
\emph{First factor in the vector variables.} Fix the covector list
$\boldsymbol\lambda=(\lambda_1,\ldots,\lambda_k)$. The map
$(v_1,\ldots,v_k)\mapsto\det(\lambda_i(v_j))$ is multilinear by the
column laws of the determinant and vanishes when two vector entries agree.
It therefore induces $L_{\boldsymbol\lambda}\in(\Lambda^kV)^*$.

\emph{Then factor in the covector variables.} Regard
$\boldsymbol\lambda\mapsto L_{\boldsymbol\lambda}$ as a map
$(V^*)^k\to(\Lambda^kV)^*$. For $\xi=v_1\wedge\cdots\wedge v_k$, the determinant row laws give
\[
\begin{aligned}
L_{(\ldots,a\lambda_j+b\mu_j,\ldots)}(\xi)
&=aL_{(\ldots,\lambda_j,\ldots)}(\xi)
+bL_{(\ldots,\mu_j,\ldots)}(\xi).
\end{aligned}
\]
Equal covector entries give equal rows, hence zero on every such $\xi$.
The resulting linear maps already exist. Their equality on those spanning
wedges proves these identities as equalities of linear maps. Thus this
map is alternating multilinear, and the exterior universal property gives
the linear $P$ with the asserted value.

\emph{Construct the inverse.} Choose a basis $(e_i)$ of $V$, with dual
$(e^i)$. Let $e_I=e_{i_1}\wedge\cdots\wedge e_{i_k}$ and
$e^I=e^{i_1}\wedge\cdots\wedge e^{i_k}$ for increasing $I$.
The determinant pairing gives $P(e^I)(e_J)=\delta_{IJ}$: matching lists
give the identity matrix; for distinct lists choose $j_s\in J\setminus I$,
so $e^{i_r}(e_{j_s})=0$ for every $r$ and column $s$ is zero, exactly
as in \cref{thm:appendix-exterior-basis}.
Define
\[
Q:(\Lambda^kV)^*\to\Lambda^k(V^*),\qquad Q(\ell)=\sum_I\ell(e_I)e^I.
\]
This finite sum is linear in $\ell$. For every basis wedge $e_J$,
\[
P(Q(\ell))(e_J)=\sum_I\ell(e_I)\delta_{IJ}=\ell(e_J).
\]
Thus $PQ(\ell)=\ell$ by equality of linear maps on a basis.
In the reverse direction
\[
Q(P(e^I))=\sum_JP(e^I)(e_J)e^J=e^I.
\]
Both sides are linear in the covector wedge, and the $e^I$ form a basis,
so $QP=\operatorname{id}$. The basis writes the inverse explicitly, but
the inverse is independent of the basis because it is the unique inverse
of the canonical $P$. In degree zero the pairing is multiplication of
scalars; if $k>n$ both spaces are zero and the empty sums give their unique maps.
\end{proof}
Combining the two identifications recovers the alternating covectors called
$\Lambda^k(V^*)$ in Chapter~13. A vector wedge is made from vectors; a
covector wedge evaluates on them and returns a scalar. For example
$(\dd x\wedge\dd z)((1,2,3),(0,1,4))=4$. The determinant pairing
relates these roles without confusing them or dividing by a factorial.

\begin{proposition}[Duality commutes with maps]
For linear $T:V\to W$ between finite-dimensional spaces, the determinant
pairings satisfy
\[
P_V\circ\Lambda^k(T^*)=(\Lambda^kT)^*\circ P_W.
\]
\end{proposition}
\begin{proof}
All four maps have been constructed. For an elementary covector wedge
$\lambda_1\wedge\cdots\wedge\lambda_k$ and an elementary vector wedge
$v_1\wedge\cdots\wedge v_k$, both sides evaluate to
$\det(\lambda_i(Tv_j))$. Spanning first in vectors, then in covectors
proves equality of the linear maps. In degree zero both sides are the identity.
\end{proof}
In words: pull covectors back before pairing, or push vectors forward before
pairing; the scalar result is the same.
Over $\R$ with $n>0$, a nonzero top covector $\omega$ is a volume form.
It selects the orientation of bases on which its value is positive.
Replacing it by a positive multiple preserves that orientation, whereas
changing its magnitude changes the normalization of volume. Evaluation
on a basis, now for an endomorphism $T:V\to V$, and the preceding determinant computation give
$T^*\omega=(\det T)\omega$. This joins the constructions: a linear map
acts on vectors, its top exterior map records the signed scaling, and a
top covector evaluates that scaling. The scalar identity works over every
field; the language of positive orientation uses real order.

\begin{proposition}[Compatibility with the shuffle product]\label{prop:abstract-shuffle-compatibility}
Let $V$ be finite dimensional over $\R$. Denote by
\[
J_k:\Lambda^k(V^*)_{\mathrm{abstract}}\to\operatorname{Alt}^k(V;\R)
\]
the determinant-pairing identification obtained above. For degrees $p,q$,
\[
J_{p+q}(\alpha\wedge\beta)
=J_p(\alpha)\wedge_{\mathrm{shuffle}}J_q(\beta).
\]
\end{proposition}
\begin{proof}
Both sides are existing bilinear maps in $\alpha,\beta$: the abstract
product was constructed in ``\nameref{sec:exterior-multiplication},'' the maps
$J_k$ by duality, and the shuffle product in \cref{sec:exterior-covectors}.
It suffices by \cref{lem:multilinear-spanning} to take
$\alpha=\lambda_1\wedge\cdots\wedge\lambda_p$ and
$\beta=\mu_1\wedge\cdots\wedge\mu_q$.
Evaluate on $v_1,\ldots,v_{p+q}$. The left side is the determinant whose
rows are the values of $\lambda_1,\ldots,\lambda_p,\mu_1,\ldots,\mu_q$
on these vectors. Expand by the columns assigned to the first $p$ rows.
For each $(p,q)$-shuffle $\sigma$ the contribution is
\[
\operatorname{sgn}(\sigma)
\det(\lambda_i(v_{\sigma(j)}))_{i,j=1}^{p}
\det(\mu_i(v_{\sigma(p+j)}))_{i,j=1}^{q}.
\]
Summing gives precisely the shuffle definition on the right. This is the
determinant/concatenation calculation already explained in Chapter~13.
Equality on every vector tuple proves equality of these alternating forms;
spanning now extends it to all $\alpha,\beta$. Degree zero uses scalar
multiplication and the empty determinant $1$; degrees above $\dim V$ give zero.
\end{proof}
Appendix~\ref{app:tensors} constructs an exterior algebra through quotients, whereas
Chapter~13 constructs it through alternating covectors and shuffles.
The canonical maps $J_k$ identify their products as well as their vector
spaces, so they give the same exterior algebra after this identification.

\section*{Exercises}
\addcontentsline{toc}{section}{Exercises}
\begin{hexercise}{app-d-1}
Let $R=\operatorname{span}(1,1)$ in $F^2$. Determine which of $x+y$ and
$x-y$ factor through $F^2/R$, including characteristic two. Construct an
isomorphism of the quotient with $F$, its reverse, and both composites.
\end{hexercise}
\begin{hexercise}{app-d-2}
Let $R\subseteq E$, $S\subseteq H$ be subspaces and $L:E\to H$ linear
with $L(R)\subseteq S$. Construct $\widetilde L:E/R\to H/S$ by using the
quotient universal property on $\pi_HL$. Verify the commutative square and
prove uniqueness. Explain why the hypothesis is necessary.
\end{hexercise}
\begin{exercise}
Construct the map $F^2\otimes F^2\to F^2$ from
$b(x,y)=(x_1y_2,x_2y_1)$, checking both slots first. Evaluate it on
$(e_1-e_2)\otimes(e_1+2e_2)$ using two different expansions.
\end{exercise}
\begin{exercise}
A proposed real-valued rule on $\R^2\otimes\R^2$ sends
$v\otimes w$ to $\|v\|_2\|w\|_2$. Compare the equal tensors
$(-v)\otimes w=v\otimes(-w)$ and the additivity requirement on
$(v+(-v))\otimes w$. Decide whether the proposed values can be those of
a linear functional, and identify the failed slot-linearity condition.
\end{exercise}
\begin{hexercise}{app-d-5}
Suppose an extra relation $v\otimes w=w\otimes v$ is imposed on
$F^2\otimes F^2$. Use $b(v,w)=v_1w_2$ to show the resulting quotient
does not satisfy the ordinary tensor universal property.
\end{hexercise}
\begin{exercise}
For $V_1=F$, $V_2=F^2$, and $W=F^2$, construct both directions of
$(V_1\oplus V_2)\otimes W\cong(V_1\otimes W)\oplus(V_2\otimes W)$.
Evaluate them on $((1,(1,-1))\otimes(2,1))$ and on a sum of two
injected elementary tensors. Verify both composites on these inputs.
\end{exercise}
\begin{hexercise}{app-d-7}
In \cref{ex:four-tensor-components}, consider the component tuple
$(0,e_1\otimes f_1+e_2\otimes f_2,0,0)$. Send it backward explicitly
and show that the result is not an elementary tensor of $A\otimes B$.
Use coordinate projections to isolate a two-by-two coefficient table and
its nonzero determinant. Explain why a dimension count alone cannot detect this.
\end{hexercise}
\begin{exercise}
For $A:V\to\mathcal L(W,U)$, distinguish the types of $A$, $A(v)$, and
$A(v)(w)$. Reconstruct its tensor factor and prove the two currying inverse identities.
\end{exercise}
\begin{exercise}
Expand $(e_1+e_2)\wedge(e_2+e_3)$ in $F^3$, and recover the three
coefficients with the induced determinant maps constructed from independent inputs.
Compute its wedge with $e_1+e_3$. Interpret the final coefficient $2$
 in characteristic two and relate its vanishing to dependence of the three vectors.
\end{exercise}
\begin{exercise}
For $T(e_1)=e_1-e_3$, $T(e_2)=2e_2$, $T(e_3)=e_1+e_3$ over $\R$,
compute $\Lambda^2T$ by its basis images and compare $\Lambda^3T$ with
$\det T$. Verify the quotient compatibility square on an elementary tensor.
Also compute $(1+e_1+e_1\wedge e_2)\wedge(2+e_2)$ in $\Lambda\R^2$.
\end{exercise}
\begin{hexercise}{app-d-11}
In the construction of exterior multiplication, prove explicitly that the
map $\mathbf v\mapsto L_{\mathbf v}$ takes equal $v$ entries to the zero
linear map. State where spanning, existence, and alternation are used.
\end{hexercise}
\begin{exercise}
For $\alpha=e^1\wedge e^3$ and $\xi=(e_1+e_2)\wedge(e_2+2e_3)$,
compute the determinant pairing and identify the spaces containing $\alpha$,
$\xi$, and the answer. Construct the inverse duality map for this basis.
\end{exercise}
